% ==========================================================
%  MASTER SUMMARY TEMPLATE — AI / ML / MATHEMATICS
%  Design Philosophy:
%    1. Readable running text is the default. Prose over boxes.
%    2. Boxes are used sparingly as visual landmarks for skimming.
%    3. Visual hierarchy: chapter > section > subsection > prose.
%    4. Every box is breakable — no orphaned whitespace at page ends.
%    5. A single consistent accent color ties all structure together.
% ==========================================================

\documentclass[11pt, oneside]{book}


% ══════════════════════════════════════════════════════════
%  PAGE GEOMETRY
% ══════════════════════════════════════════════════════════
\usepackage[
  a4paper,
  top    = 2.6cm,
  bottom = 2.1cm,
  left   = 2.5cm,
  right  = 2.5cm,
  headheight = 15pt
]{geometry}


% ══════════════════════════════════════════════════════════
%  TYPOGRAPHY
% ══════════════════════════════════════════════════════════
\usepackage[T1]{fontenc}
\usepackage[utf8]{inputenc}
\usepackage{lmodern}         % Latin Modern: improved Computer Modern
\usepackage{microtype}       % Microtypographic refinements (better justification,
                             %   fewer overfull hboxes — invaluable for math text)

% If inconsolata is not installed, simply remove this line.
% It provides a crisper monospace font for code listings.
\usepackage{inconsolata}

\setlength{\parindent}{0pt}
\setlength{\parskip}{0.62em}


% ══════════════════════════════════════════════════════════
%  MATHEMATICS
% ══════════════════════════════════════════════════════════
\usepackage{amsmath, amssymb, amsthm, mathtools}
\usepackage{bm}        % Bold math symbols: \bm{\theta}
\usepackage{bbm}       % Blackboard bold for indicator: \mathbbm{1}
\usepackage{mathrsfs}  % Script letters: \mathscr{F}, \mathscr{H}
\usepackage{nicefrac}  % Compact inline fractions: \nicefrac{1}{2}
\usepackage{cancel}    % Strike-through in math: \cancel{x}


% ══════════════════════════════════════════════════════════
%  SCIENTIFIC NOTATION
% ══════════════════════════════════════════════════════════
% siunitx — consistent typesetting of numbers and SI units.
%   Usage: \SI{1.5}{\micro\meter}  \SI{50}{\kilo\hertz}  \num{6.022e23}
\usepackage{siunitx}
\sisetup{per-mode = symbol}   % render "per unit" as e.g. m/s rather than m s^{-1}

% mhchem — chemical formulas and reaction equations.
%   Usage: \ce{H2O}   \ce{ATP -> ADP + P_i}   \ce{^{14}C}
\usepackage[version=4]{mhchem}


% ══════════════════════════════════════════════════════════
%  FIGURES, TABLES, LISTS
% ══════════════════════════════════════════════════════════
\usepackage{graphicx}
\usepackage{booktabs}         % \toprule, \midrule, \bottomrule
\usepackage{array}            % Extended column types in tabular
\usepackage{float}            % [H] float placement
\usepackage{enumitem}

\setlist[itemize]{topsep=3pt, itemsep=2pt, leftmargin=2.2em}
\setlist[enumerate]{topsep=3pt, itemsep=2pt, leftmargin=2.4em}

\usepackage{needspace}        % keep headings together with their first lines

% ── Course addition: self-drawn figures (TikZ / pgfplots) ──
\usepackage{tikz}
\usetikzlibrary{arrows.meta, positioning, shapes.geometric, calc, fit, backgrounds}
\usepackage{pgfplots}
\pgfplotsset{compat=1.18}


% ══════════════════════════════════════════════════════════
%  ALGORITHMS / PSEUDOCODE
% ══════════════════════════════════════════════════════════
% Remove this block if you do not need pseudocode.
\usepackage[ruled, vlined, linesnumbered]{algorithm2e}
\SetAlFnt{\small}
\SetAlCapFnt{\small\bfseries}
\SetCommentSty{textit}


% ══════════════════════════════════════════════════════════
%  COLORS  (a refined, academic palette)
% ══════════════════════════════════════════════════════════
\usepackage{xcolor}

% Primary structural accent — deep navy used for headings, rules, ToC
\definecolor{accent}{HTML}{1B3A5C}

% Box accent colors: frame and background, designed to be legible but
% unobtrusive so running text remains the visual focus.
\definecolor{def@frame}{HTML}{1A5276}   \definecolor{def@bg}{HTML}{EBF5FB}
\definecolor{thm@frame}{HTML}{922B21}   \definecolor{thm@bg}{HTML}{FDEDEC}
\definecolor{ex@frame} {HTML}{1E8449}   \definecolor{ex@bg} {HTML}{EAFAF1}
\definecolor{int@frame}{HTML}{9A6A00}   \definecolor{int@bg}{HTML}{FEF9E7}
\definecolor{note@frame}{HTML}{5B2C6F}  \definecolor{note@bg}{HTML}{F5EEF8}


% ══════════════════════════════════════════════════════════
%  CODE LISTINGS
% ══════════════════════════════════════════════════════════
\usepackage{listings}
\lstset{
  basicstyle       = \ttfamily\small,
  breaklines       = true,
  numbers          = left,
  numberstyle      = \tiny\color{gray!70},
  frame            = tb,
  framerule        = 0.4pt,
  rulecolor        = \color{gray!35},
  tabsize          = 2,
  showstringspaces = false,
  keywordstyle     = \color{def@frame}\bfseries,
  commentstyle     = \color{gray!55!black}\itshape,
  stringstyle      = \color{ex@frame!90!black},
  backgroundcolor  = \color{gray!5},
  xleftmargin      = 2em,
  xrightmargin     = 0.4em,
  aboveskip        = 0.8em,
  belowskip        = 0.5em
}


% ══════════════════════════════════════════════════════════
%  CUSTOM BOXES  (tcolorbox)
%
%  KEY DESIGN DECISIONS:
%  1. enhanced   — enables advanced tcolorbox features (borderline etc.)
%  2. breakable  — box may split across a page break; eliminates the ugly
%                  white gap that appears when a long box would overflow.
%  3. Left-accent bar style — a thick colored stripe on the left margin.
%                  This is legible, modern, and unobtrusive relative to
%                  a full colored frame. The box background is a very
%                  light tint, keeping the reading experience calm.
%  4. sharp corners=west — left side is flat (squared), right is rounded.
%                  This mirrors the visual logic of the left accent bar.
%  5. parbox=false — allows display math to stretch properly inside boxes.
% ══════════════════════════════════════════════════════════
\usepackage[most]{tcolorbox}

\tcbset{
  enhanced,
  breakable,
  arc             = 2.5mm,
  boxrule         = 0.5pt,
  left            = 4.5mm,
  right           = 4mm,
  top             = 2.5mm,
  bottom          = 2.5mm,
  fonttitle       = \bfseries\small,
  sharp corners   = west,
  parbox          = false
}

\newtcolorbox{defbox}[1][]{
  title  = {Definition},
  colback = def@bg, colframe = def@frame,
  borderline west = {3.5pt}{0pt}{def@frame},
  #1
}

\newtcolorbox{thmbox}[1][]{
  title  = {Theorem / Key Formula},
  colback = thm@bg, colframe = thm@frame,
  borderline west = {3.5pt}{0pt}{thm@frame},
  #1
}

\newtcolorbox{exbox}[1][]{
  title  = {Example},
  colback = ex@bg, colframe = ex@frame,
  borderline west = {3.5pt}{0pt}{ex@frame},
  #1
}

\newtcolorbox{intbox}[1][]{
  title  = {Intuition \& Mental Model},
  colback = int@bg, colframe = int@frame,
  borderline west = {3.5pt}{0pt}{int@frame},
  #1
}

\newtcolorbox{notebox}[1][]{
  title  = {Note},
  colback = note@bg, colframe = note@frame,
  borderline west = {3.5pt}{0pt}{note@frame},
  #1
}


% ══════════════════════════════════════════════════════════
%  TABLE OF CONTENTS  (tocloft)
%
%  The [titles] option prevents tocloft from conflicting
%  with titlesec's redefinition of chapter/section titles.
% ══════════════════════════════════════════════════════════
\usepackage[titles]{tocloft}

% Chapter entries: accented, bold, with dots leading to page number
\renewcommand{\cftchapfont}     {\bfseries\color{accent}}
\renewcommand{\cftchappagefont} {\bfseries\color{accent}}
\renewcommand{\cftchapleader}   {\bfseries\color{accent!40}\cftdotfill{\cftdotsep}}
\setlength{\cftbeforechapskip} {7pt}

% Section entries: standard weight, small size
\renewcommand{\cftsecfont}      {\small}
\renewcommand{\cftsecpagefont}  {\small}
\setlength{\cftsecindent}      {1.8em}
\setlength{\cftbeforesecskip}  {1.5pt}

% Subsection entries: muted gray to visually recede
\renewcommand{\cftsubsecfont}      {\small\color{gray!65!black}}
\renewcommand{\cftsubsecpagefont}  {\small\color{gray!65!black}}
\setlength{\cftsubsecindent}      {3.4em}
\setlength{\cftbeforesubsecskip}  {0.5pt}


% ══════════════════════════════════════════════════════════
%  CHAPTER & SECTION HEADINGS  (titlesec)
%
%  Chapter: large faded chapter number floated right, title
%           left-aligned beneath it, then a double rule.
%           The oversized ghost number is a quiet nod to
%           modern textbook design without being gimmicky.
%
%  Section: accent-colored, with a thin rule beneath.
%           The rule draws the eye without shouting.
%
%  Subsection: slightly smaller accent color, no rule — the
%           rule at section level is already enough structure.
% ══════════════════════════════════════════════════════════
\usepackage{titlesec}

\titleformat{\chapter}[display]
  {\normalfont\bfseries\color{accent}\huge}
  {\hfill\fontsize{68}{68}\selectfont\color{accent!13}\thechapter}
  {-10pt}
  {}
  [{\vspace{4pt}%
    {\color{accent}\rule{\linewidth}{1.6pt}}\\[-2pt]%
    {\color{accent!35}\rule{\linewidth}{0.5pt}}}]
\titlespacing*{\chapter}{0pt}{-28pt}{22pt}

\titleformat{\section}
  {\normalfont\large\bfseries\color{accent}}
  {\thesection}
  {0.8em}
  {}
  [{\vspace{1pt}\color{accent!45}\titlerule[0.45pt]}]
\titlespacing*{\section}{0pt}{3.8ex plus 1ex minus .2ex}{2.2ex plus .2ex}

\titleformat{\subsection}
  {\normalfont\normalsize\bfseries\color{accent!80!black}}
  {\thesubsection}
  {0.7em}
  {}
\titlespacing*{\subsection}{0pt}{2.6ex plus 1ex minus .2ex}{1.4ex plus .2ex}

% Run-in paragraph label — useful for named sub-items in a derivation.
\titleformat{\paragraph}[runin]
  {\normalfont\small\bfseries}
  {}
  {0em}
  {}[.\enspace]
\titlespacing{\paragraph}{0pt}{1.2ex plus .5ex minus .2ex}{0.5em}


% ══════════════════════════════════════════════════════════
%  HEADER & FOOTER
% ══════════════════════════════════════════════════════════
\usepackage{fancyhdr}
\pagestyle{fancy}
\fancyhf{}
\fancyhead[L]{\small\color{gray!70}\nouppercase{\leftmark}}
\fancyhead[R]{\small\color{gray!70}\thepage}
\renewcommand{\headrulewidth}{0.4pt}

% Chapter-opening pages use LaTeX's 'plain' style — redefine it here
% to match the rest of the document (page number at foot, no header line).
\fancypagestyle{plain}{%
  \fancyhf{}
  \fancyfoot[C]{\small\color{gray!70}\thepage}
  \renewcommand{\headrulewidth}{0pt}%
}


% ══════════════════════════════════════════════════════════
%  LINKS & REFERENCES
% ══════════════════════════════════════════════════════════
\usepackage[hidelinks]{hyperref}
\usepackage[nameinlink, capitalize]{cleveref}


% ══════════════════════════════════════════════════════════
%  DOCUMENT METADATA  (fill in for each new subject)
% ══════════════════════════════════════════════════════════
\newcommand{\coursetitle}{Probabilistic Models}
\newcommand{\documenttype}{Study Notes -- Working Draft}
\newcommand{\authorname}{Stefan Eder}
\newcommand{\basedon}{Based on the lecture \emph{Probabilistic Models} (G.~Widmer, JKU Linz), WS 2026/27}


% ══════════════════════════════════════════════════════════
%  MATHEMATICAL MACROS
%  Organized by domain. Add your own at the end of each group.
% ══════════════════════════════════════════════════════════

% ── Number Sets & Spaces ──────────────────────────────────
\newcommand{\R}{\mathbb{R}}
\newcommand{\N}{\mathbb{N}}
\newcommand{\Z}{\mathbb{Z}}
\newcommand{\C}{\mathbb{C}}
\newcommand{\Q}{\mathbb{Q}}

% ── Probability & Statistics ──────────────────────────────
\newcommand{\E}{\mathbb{E}}              % Expectation
\newcommand{\Prob}{\mathbb{P}}           % Probability measure
\newcommand{\Var}{\mathrm{Var}}
\newcommand{\Cov}{\mathrm{Cov}}
\newcommand{\Corr}{\mathrm{Corr}}
\newcommand{\iid}{\overset{\mathrm{iid}}{\sim}}
% Indicator function. Usage: \ind_{\{x > 0\}}
\newcommand{\ind}{\mathbbm{1}}

% ── Information Theory ────────────────────────────────────
% KL divergence. Usage: \KL{q}{p}  →  D_KL(q ‖ p)
\newcommand{\KL}[2]{D_{\mathrm{KL}}\!\left(#1 \,\|\, #2\right)}
\newcommand{\Ent}{\mathrm{H}}            % Entropy:         \Ent(X)
\newcommand{\MI}{\mathrm{I}}             % Mutual information: \MI(X;Y)

% ── Calculus & Optimization ───────────────────────────────
\newcommand{\argmax}{\operatorname*{arg\,max}}
\newcommand{\argmin}{\operatorname*{arg\,min}}
\newcommand{\grad}{\nabla}               % Gradient
\newcommand{\hess}{\nabla^2}             % Hessian
\newcommand{\lapl}{\Delta}               % Laplacian
% Differential (upright d). Usage: \int f(x) \dd x
\newcommand{\dd}{\,\mathrm{d}}
% Partial derivatives. Usage: \pd{f}{x}  or  \pdd{f}{x}
\newcommand{\pd}[2]{\frac{\partial #1}{\partial #2}}
\newcommand{\pdd}[2]{\frac{\partial^2 #1}{\partial #2^2}}

% ── Norms, Abs, Inner Products ────────────────────────────
\newcommand{\norm}[1]{\left\lVert #1 \right\rVert}
\newcommand{\normsq}[1]{\left\lVert #1 \right\rVert^2}
\newcommand{\abs}[1]{\left\lvert #1 \right\rvert}
\newcommand{\inner}[2]{\left\langle #1,\, #2 \right\rangle}

% ── Linear Algebra ────────────────────────────────────────
\newcommand{\T}{^{\mathsf{T}}}           % Transpose:        \mathbf{A}\T
\newcommand{\inv}{^{-1}}                 % Inverse:          A\inv
\newcommand{\pinv}{^{\dagger}}           % Pseudo-inverse:   A\pinv
\newcommand{\tr}{\operatorname{tr}}
\newcommand{\rank}{\operatorname{rank}}
\newcommand{\diag}{\operatorname{diag}}
\newcommand{\spanop}{\operatorname{span}}
\newcommand{\eye}{\mathbf{I}}            % Identity matrix
\newcommand{\zeros}{\mathbf{0}}
\newcommand{\ones}{\mathbf{1}}

% ── ML / AI Core Notation ────────────────────────────────
\newcommand{\loss}{\mathcal{L}}          % Loss:          \loss(\params)
\newcommand{\reg}{\mathcal{R}}           % Regularizer:   \reg(\params)
\newcommand{\dataset}{\mathcal{D}}       % Dataset:       \dataset
\newcommand{\params}{\boldsymbol{\theta}}   % Parameters
\newcommand{\weights}{\mathbf{w}}
\newcommand{\xvec}{\mathbf{x}}
\newcommand{\yvec}{\mathbf{y}}
\newcommand{\zvec}{\mathbf{z}}
\newcommand{\hvec}{\mathbf{h}}           % Hidden state
\newcommand{\feature}{\boldsymbol{\phi}} % Feature map

% Activation functions
\DeclareMathOperator{\softmax}{softmax}
\DeclareMathOperator{\relu}{ReLU}
\DeclareMathOperator{\sign}{sign}
% sigmoid: just use \sigma (standard Greek letter)

% ── Complexity ────────────────────────────────────────────
\newcommand{\bigO}[1]{\mathcal{O}\!\left(#1\right)}
\newcommand{\bigOm}[1]{\Omega\!\left(#1\right)}
\newcommand{\bigTh}[1]{\Theta\!\left(#1\right)}

% ── Common Distributions ──────────────────────────────────
\newcommand{\Normal}{\mathcal{N}}
\newcommand{\Uniform}{\mathcal{U}}
\newcommand{\Bernoulli}{\mathrm{Ber}}
\newcommand{\Categorical}{\mathrm{Cat}}
\newcommand{\Dirichlet}{\mathrm{Dir}}
\newcommand{\Binomial}{\mathrm{Bin}}

% ── Miscellaneous ─────────────────────────────────────────
\newcommand{\given}{\,\vert\,}           % Conditional: \Prob(A \given B)
\newcommand{\st}{\;\text{s.t.}\;}        % "subject to"
\newcommand{\eqdef}{\coloneqq}           % Definitional equality  :=
\newcommand{\eqq}{\overset{?}{=}}        % "Does this equal?"



% ══════════════════════════════════════════════════════════
%  COURSE-SPECIFIC MACROS  (Probabilistic Models)
%  Notation follows the lecture (Koller & Friedman style):
%  calligraphic X = all variables, bold = sets, superscripts = values.
% ══════════════════════════════════════════════════════════
\newcommand{\Val}{\mathit{Val}}          % domain:  \Val(X)
\newcommand{\Pa}{\mathrm{Pa}}            % parents: \Pa(X_i)
\newcommand{\NonDesc}{\mathrm{NonDesc}}  % non-descendants
\newcommand{\Desc}{\mathrm{Desc}}        % descendants
\newcommand{\indep}{\perp}               % (X \indep Y \given Z)
\newcommand{\nindep}{\not\perp}          % not independent
\newcommand{\Xall}{\mathcal{X}}          % complete set of variables
\newcommand{\bX}{\mathbf{X}}  \newcommand{\bY}{\mathbf{Y}}
\newcommand{\bZ}{\mathbf{Z}}  \newcommand{\bE}{\mathbf{E}}
\newcommand{\bx}{\mathbf{x}}  \newcommand{\by}{\mathbf{y}}
\newcommand{\bz}{\mathbf{z}}  \newcommand{\be}{\mathbf{e}}
\newcommand{\vn}[1]{\mathit{#1}}         % named variable/value: \vn{Flu}, \vn{flu}

% ── Figure colours (aliases of the template palette) ──────
\colorlet{figblue}{def@frame}
\colorlet{figred}{thm@frame}
\colorlet{figgreen}{ex@frame}
\colorlet{figamber}{int@frame}
\colorlet{figpurple}{note@frame}

% ── TikZ styles for Bayesian-network diagrams ─────────────
\tikzset{
  bn/.style     = {ellipse, draw=accent, thick, fill=accent!6,
                   minimum height=7mm, minimum width=17mm,
                   inner sep=1.5pt, font=\small\itshape},
  bnsmall/.style= {circle, draw=accent, thick, fill=accent!6,
                   minimum size=7.5mm, inner sep=0pt, font=\small\itshape},
  edge/.style   = {-{Stealth[length=2.3mm]}, thick, accent!85!black},
  evid/.style   = {fill=gray!30},
  qry/.style    = {draw=figred, very thick, fill=figred!8},
  cpd/.style    = {font=\scriptsize, inner sep=1pt, fill=white},
  cpdlink/.style= {gray!55, thin},
}

% ── Placeholders ──────────────────────────────────────────
% Figure placeholder (same signature as in other course notes).
% #1 height, #2 concept, #3 what the picture should show + search hint,
% #4 caption, #5 label.  Replace with \includegraphics{figures/...}
% once a suitable image is in figures/.
\newcommand{\figureplaceholder}[5]{%
\begin{figure}[htbp]
\centering
\begin{tcolorbox}[unbreakable,width=\linewidth,height=#1,colback=accent!3,
  colframe=accent!30,boxrule=0.5pt,arc=2.5mm,sharp corners=all,
  valign=center,halign=center,fontupper=\small]
{\bfseries\color{accent}Figure to insert: #2}\par\medskip
{\color{gray!75!black}#3}
\end{tcolorbox}
\caption{#4}\label{#5}
\end{figure}}

% Outline text for chapters that are not written yet.
\newcommand{\planned}[1]{%
  \par{\small\color{gray!70!black}\textit{Planned content:} #1}\par}

% Keep section headings together with the start of their text.
\let\originalsection\section
\renewcommand{\section}{\Needspace{7\baselineskip}\originalsection}
\let\originalsubsection\subsection
\renewcommand{\subsection}{\Needspace{6\baselineskip}\originalsubsection}

% ==========================================================
%  BEGIN DOCUMENT
% ==========================================================
\begin{document}

\frontmatter  % Roman-numeral page numbers for front matter


% ──────────────────────────────────────────────────────────
%  TITLE PAGE
% ──────────────────────────────────────────────────────────
\thispagestyle{empty}
\vspace*{1.8cm}

\begin{center}
  {\color{accent}\rule{\linewidth}{1.8pt}}
  \vspace{0.65cm}

  {\fontsize{28}{34}\selectfont\bfseries\color{accent}\coursetitle\par}
  \vspace{0.25cm}
  {\large\color{gray!80!black}\documenttype\par}

  \vspace{0.55cm}
  {\color{accent!35}\rule{\linewidth}{0.5pt}}
  \vspace{0.55cm}

  {\large\authorname\par}
  \vspace{0.2cm}
  {\small\color{gray!70}\basedon\par}
\end{center}

\vspace{1.6cm}

\begin{center}
\begin{minipage}{0.84\textwidth}
  \small\color{gray!80!black}
  These notes try to make probabilistic models \emph{click}. Every idea is
  introduced in plain words first, then as a formula, and whenever things get
  abstract there is a small example with actual numbers. All worked numbers
  were re-computed independently.
  Use the notes as a study aid alongside the lecture slides, not as a
  replacement for them.

  \medskip
  \textbf{Working draft:} Chapters 1--3 are written out. Chapters 4--8 are
  still outlines and will be filled in as the lecture progresses.

  \medskip
  If you spot an error or a gap, please open an issue or pull request at
  \textbf{github.com/steveTgit/AI\_Master} --- corrections are always welcome.

  \medskip
  \color{gray!60!black}
  Good study materials should be shared, not gatekept.
  If this helped you, take it as a reason to help the next colleague you see struggling with something.
\end{minipage}
\end{center}

\vfill
{\centering\small\color{gray!65} \today\par}
\clearpage


% ──────────────────────────────────────────────────────────
%  TABLE OF CONTENTS
% ──────────────────────────────────────────────────────────
\setcounter{tocdepth}{2}  % Show: chapters, sections, subsections only

\begingroup
  \let\cleardoublepage\relax
  \let\clearpage\relax
  \tableofcontents
\endgroup

\mainmatter  % Switch to Arabic numerals




% ==========================================================
%  CHAPTER 1 — INTRODUCTION            (lecture part 1, VO 1)
% ==========================================================
\chapter{Introduction: Why Probabilistic Models?}
\label{ch:intro}

Almost everything interesting that an intelligent system does involves
information that is incomplete, noisy, or simply wrong. A sensor flickers, a
symptom has five possible causes, a microphone picks up the audience coughing.
This course is about one systematic answer to the question \emph{how should a
machine reason when it cannot be sure?} The answer: describe the world with
probabilities, organise those probabilities with a graph, and let
general-purpose algorithms do the reasoning.

This first chapter sets the scene. We look at the kind of problems we want to
solve, why logic alone does not get us there, and which three big problems the
rest of the course is going to tackle.

\section{Reasoning Under Uncertainty}

The lecture opens with a handful of scenarios that look very different on the
surface:

\begin{itemize}
  \item \textbf{Diagnosing a complex system.} A spacecraft (or one of its
    subsystems) sends sensor signals that indicate errors --- ``Houston, we
    have a problem''. Which fault most likely caused them? What would happen if
    we tried a certain repair? Which repair should we actually do?
  \item \textbf{Autonomous driving.} Given a road map and streams of GPS,
    camera, radar and range-sensor data, drive a car through difficult terrain,
    with other cars and obstacles around.
  \item \textbf{Aircraft tracking.} From intermittent, sometimes occluded radar
    or video signals: where will each aircraft be in 0.1, 1 or 10 seconds, how
    many aircraft are there, and which one is which?
  \item \textbf{Speech recognition.} Split an audio stream into words,
    transcribe them, and use spoken commands to control devices.
  \item \textbf{Real-time music following.} Listen to a live performance and
    keep track of where the musicians currently are in the printed score ---
    for example to turn pages automatically or to synchronise other events with
    the music.
\end{itemize}

Underneath, all of them have the same structure. An autonomous system receives
information about its environment, at regular or irregular intervals. From
that information it has to identify objects, classify situations, make
predictions and take decisions. And the information is incomplete, uncertain,
noisy, sometimes wrong and occasionally contradictory.

That last point is what breaks purely logical or deterministic approaches. Try
writing down the rule ``muscle pain $\Rightarrow$ flu''. It is wrong for most
people with muscle pain. Patch it with exceptions (``\ldots unless they went to
the gym, unless they slept badly, unless \ldots'') and the list never ends.
What we actually know is softer: \emph{muscle pain makes flu more likely}, and
we would like to say \emph{how much} more likely. That is exactly what
probabilities are for.

Probabilistic graphical models give us four things at once. They
\begin{itemize}
  \item represent knowledge about the world \emph{and} about its uncertainty,
  \item support inference: drawing conclusions and making decisions,
  \item can be learned from example observations, and
  \item technically, are graphical representations of \emph{factorised}
    probability distributions.
\end{itemize}
The last point will only make sense after \cref{ch:bn}; by then it should feel
obvious. The field sits where artificial intelligence, statistics, probability
theory, machine learning and algorithms meet.

\section{The Declarative Approach}

There are two ways to build a system that solves problems. The
\emph{procedural} way is to write a special-purpose algorithm for each problem.
The \emph{declarative} way --- the classic AI philosophy, and the one this
course follows --- is to
\begin{enumerate}
  \item build a \emph{model} of the problem or world you want to reason about,
    and
  \item write \emph{general inference algorithms} that work on any such model
    and draw conclusions from it.
\end{enumerate}

This buys you a lot. Knowledge (the model) and reasoning (the algorithm) are
cleanly separated. The model has a clear, well-defined meaning, independent of
whatever algorithm happens to run on it. The inference algorithms are generic:
they do not care which domain the model describes. And if the world changes,
you only change the model, not the code.

\begin{intbox}
Think of a city map versus a set of written directions. Directions
(procedural) get you along \emph{one} route; if a street is closed, they are
useless. A map (declarative) answers \emph{any} routing question with one
general skill --- reading a map --- and if a street closes, you update the map,
not every set of directions. A probabilistic model is the map; inference
algorithms are the map-reading skill.
\end{intbox}

\section{What Are Probabilistic Graphical Models?}

A probabilistic graphical model is a formalism for describing the structure of
a domain and for quantifying how uncertain we are about it. It is called
\emph{probabilistic} because it represents a probability distribution, and
\emph{graphical} because the dependencies between the variables are made
explicit in a graph.

Three questions run through the whole course:

\paragraph{Representation (modelling)} How do we represent the relevant
aspects of a world and the connections between them?

\paragraph{Inference} Given incomplete, uncertain, noisy observations and the
model, how do we reach conclusions about things we cannot observe, or about the
future?

\paragraph{Learning} How do we build models of complex worlds automatically
from (lots of) observations?

\Cref{tab:roadmap} shows how the chapters of these notes map to the lecture
units. The course covers \emph{directed} models only: static Bayesian networks
and temporal models (dynamic Bayesian networks, hidden Markov models and the
Kalman filter). Undirected and semi-directed models (Markov networks, chain
graphs) are not part of the course.

\begin{table}[htbp]
\centering
\small
\begin{tabular}{@{}lll@{}}
\toprule
\textbf{Chapter} & \textbf{Topic} & \textbf{Lecture units} \\
\midrule
\ref{ch:intro} & Introduction & VO 1 \\
\ref{ch:prob}  & Fundamentals of probability: inference and independence & VO 1--2 \\
\ref{ch:bn}    & Bayesian networks: representation and semantics & VO 3 \\
\ref{ch:inference} & Inference in Bayesian networks: exact, approximate, MCMC & VO 4--5 \\
\ref{ch:learning}  & Learning Bayesian networks: parameters and structure & VO 6--8 \\
\ref{ch:dbn}   & Temporal models: dynamic Bayesian networks & VO 9 \\
\ref{ch:hmm}   & Hidden Markov models: inference, learning, applications & VO 10--11 \\
\ref{ch:kalman} & Continuous linear systems: the Kalman filter & video \\
\bottomrule
\end{tabular}
\caption{Roadmap of the course. Representation, inference and learning first
for static models (Chapters 3--5), then for models that evolve over time
(Chapters 6--8).}
\label{tab:roadmap}
\end{table}

\section{A First Example: A Tiny Diagnosis World}
\label{sec:diagnosis-world}

To make this concrete, consider a toy medical diagnosis setting. There are two
possible diseases, \emph{flu} and \emph{hayfever}; two possible symptoms,
(sinus) \emph{congestion} and \emph{muscle pain}; and the \emph{season}, which
might also matter. We want a system that answers questions like:
\begin{itemize}
  \item How likely is it that a patient has the flu, given that it is fall and
    she has a sinus congestion but no muscle pain?
  \item How likely is it that a patient does \emph{not} have hayfever, given
    that she has muscle pain and the flu?
  \item In which season are people most likely to have a congestion but no
    hayfever?
\end{itemize}

The probabilistic approach takes three steps. First, model every relevant
entity as a \emph{random variable}. Second, model the relationships and
influences between them as (conditional) probabilities. Third, treat question
answering as \emph{computing conditional probabilities}, given whatever has
been observed. Here, five variables are enough:
\[
\begin{array}{rcl}
  \vn{Flu}, \;\vn{Hayfever}, \;\vn{Congestion}, \;\vn{MusclePain}
    & \in & \{\vn{true}, \vn{false}\}, \\
  \vn{Season} & \in & \{\vn{spring}, \vn{summer}, \vn{fall}, \vn{winter}\},
\end{array}
\]
so the model is defined over
$\Xall = \{\vn{Flu}, \vn{Hayfever}, \vn{Congestion}, \vn{MusclePain}, \vn{Season}\}$.

Implicitly, we assume that the variables describe the state of the world (one
specific patient) at one point in time, that this state is completely
characterised by the values of the variables, that every variable always takes
exactly one of its values, and that some of these values may be known to us
while others are hidden.

The second part of the model is the \emph{full joint distribution}: one
probability for every possible combination of values. It is a table that sums
to one and here has $2 \times 2 \times 2 \times 2 \times 4 = 64$ entries. The
first few rows might look like this:
\[
\begin{array}{rcl}
  P(\neg\vn{flu}, \neg\vn{hayfever}, \neg\vn{congestion}, \neg\vn{musclepain}, \vn{spring}) &=& 0.014 \\
  P(\vn{flu}, \neg\vn{hayfever}, \neg\vn{congestion}, \neg\vn{musclepain}, \vn{spring}) &=& 0.003 \\
  P(\neg\vn{flu}, \vn{hayfever}, \neg\vn{congestion}, \neg\vn{musclepain}, \vn{spring}) &=& 0.008 \\
  \multicolumn{3}{c}{\vdots}
\end{array}
\]

The good news, shown in \cref{ch:prob}: \emph{every} question of the type above
is a conditional probability query, such as
$P(\vn{Flu} = \vn{true} \given \vn{Season} = \vn{spring}, \vn{MusclePain} = \vn{false})$,
and every such query can be answered from the full joint distribution with one
universal algorithm. The bad news comes in three flavours:

\begin{itemize}
  \item \textbf{Representational complexity.} The joint table grows
    exponentially with the number of variables. With 50 binary variables it
    has about $10^{15}$ entries --- far too many to store.
  \item \textbf{Computational complexity.} Answering a query exactly from the
    table also takes exponential time.
  \item \textbf{Modelling complexity.} Even if we could store the table, who
    would give us all those numbers? Building a valid model of a complex world
    by hand is hard.
\end{itemize}

The rest of the course is about fighting these three problems. Against the
first, we study \emph{(conditional) independence}, use it to break the joint
distribution into a product of small pieces (``factorise'' it), and draw the
result as a compact graph that shows the (in)dependencies --- and often the
causal structure --- of the domain (\cref{fig:diagnosis-net}). Against the
second, we design algorithms that compute answers directly from the graph,
optimise them, and develop fast approximate algorithms. Against the third, we
\emph{learn} models from observations, which is where machine learning comes
in. All of this needs a solid grip on basic probability theory, which is what
\cref{ch:prob} is for.

\begin{figure}[htbp]
\centering
\begin{tikzpicture}
  \node[bn] (S) at (0,0)       {Season};
  \node[bn] (F) at (-1.9,-1.4) {Flu};
  \node[bn] (H) at (1.9,-1.4)  {Hayfever};
  \node[bn, minimum width=22mm] (M) at (-3.6,-2.8) {MusclePain};
  \node[bn, minimum width=22mm] (C) at (0,-2.8)    {Congestion};
  \draw[edge] (S) -- (F);
  \draw[edge] (S) -- (H);
  \draw[edge] (F) -- (M);
  \draw[edge] (F) -- (C);
  \draw[edge] (H) -- (C);
\end{tikzpicture}
\caption{A graphical model for the diagnosis world. The season influences both
diseases; flu causes muscle pain and congestion; hayfever causes congestion.
With this structure the joint distribution needs only 17 numbers instead of 63
(we will count them in \cref{sec:compactness}).}
\label{fig:diagnosis-net}
\end{figure}


% ==========================================================
%  CHAPTER 2 — FUNDAMENTALS OF PROBABILITY   (part 2, VO 1--2)
% ==========================================================
\chapter{Fundamentals of Probability}
\label{ch:prob}

Everything later in the course is built from a surprisingly small toolbox:
random variables, joint distributions, conditioning, a couple of rules for
rewriting probabilities, and --- most importantly --- \emph{independence}. This
chapter fixes the notation used in the lecture, refreshes the basics, and then
makes two big points. First, \emph{any} probabilistic question can be answered
from the full joint distribution with one simple algorithm, inference by
enumeration. Second, that algorithm is hopelessly expensive unless we exploit
(conditional) independencies --- which is the whole motivation for graphical
models.

If you are comfortable with probability, skim the first sections, but read
\cref{sec:enumeration,sec:cond-indep} carefully: they are the foundation for
everything that follows.

% ----------------------------------------------------------
\section{Random Variables and Events}
% ----------------------------------------------------------

\begin{defbox}[title={Random variable}]
A \emph{random variable} $X$ is a variable with a fixed (finite or infinite)
domain $\Val(X)$, the set of values it can take. It represents one aspect of
the world (or of the system's internal state).
\end{defbox}

Random variables are the atomic building blocks of all our models. They come in
three flavours:
\begin{itemize}
  \item \textbf{Boolean} variables take values in $\{\vn{false}, \vn{true}\}$
    (or $\{0,1\}$).
  \item \textbf{Discrete} (categorical) variables have a finite set of symbolic
    values, e.g.\ $\Val(\vn{Season}) = \{\vn{spring}, \vn{summer}, \vn{fall},
    \vn{winter}\}$.
  \item \textbf{Continuous} variables take real values. Most real-world
    sensors (temperature, sonar, radar, \ldots) produce these; we get to them
    in \cref{sec:continuous}.
\end{itemize}

The running example of this chapter is the diagnosis world from
\cref{sec:diagnosis-world}: two diseases (not directly observable), two
symptoms (possibly observable) and the season (observable), with exactly one
patient at a time. Notice the modelling choice: each disease and each symptom
gets its \emph{own} Boolean variable. A single variable
$\vn{Disease} \in \{\vn{flu}, \vn{hayfever}, \vn{none}\}$ would be unable to
express a patient who has \emph{both} diseases.

\begin{defbox}[title={Events and atomic events}]
An \emph{event} is an assignment of values to some or all of the variables of
the world. An \emph{atomic event} is an event in which \emph{every} variable
has a value assigned.
\end{defbox}

An atomic event is one complete, fully specified state of the world. A general
event fixes only some variables and therefore stands for the \emph{set} of all
world states that agree on those values. In the diagnosis world:
\begin{itemize}
  \item $\vn{Season} = \vn{spring} \wedge \vn{MusclePain} = \vn{true}$ covers
    all situations in which it is spring and the patient has muscle pain,
    whatever her diseases or congestion;
  \item $\vn{Congestion} = \vn{false}$ is an event too;
  \item $\vn{Season} = \vn{fall} \wedge \vn{MusclePain} = \vn{true} \wedge
    \vn{Congestion} = \vn{false} \wedge \vn{Hayfever} = \vn{false} \wedge
    \vn{Flu} = \vn{false}$ is an atomic event.
\end{itemize}
Altogether there are $2 \times 2 \times 2 \times 2 \times 4 = 64$ atomic
events. Events are what the system computes probabilities for.

\begin{intbox}[unbreakable]
Picture the 64 atomic events as the 64 cells of one big spreadsheet, each cell
holding a probability. Every event is simply a \emph{selection of cells}, and
its probability is the sum over the selected cells. Marginals, conditionals and
queries later in this chapter are all just clever ways of selecting and summing
cells.
\end{intbox}

\subsection{Notation}

The lecture uses a fixed set of conventions (mostly those of Koller \&
Friedman); \cref{tab:notation} collects them. It is worth learning them early
--- the exam uses the same notation.

\begin{table}[htbp]
\centering
\small
\renewcommand{\arraystretch}{1.15}
\begin{tabular}{@{}ll@{}}
\toprule
\textbf{Example} & \textbf{Meaning} \\
\midrule
$\Xall$ & the complete set of variables a model is defined over \\
$A,\ X_i,\ E_j$ & a single random variable (e.g.\ $X_i \in \Xall$) \\
$\bX, \bY, \bZ, \bE, \ldots$ & sets of random variables, e.g.\ $\bX = \{X_1, X_2, X_3\}$ \\
$\Val(X)$ & domain of $X$: the set of values $X$ can take \\
$x^0, x^1, \ldots$ & the values of $X$; $x^i$ is the $i$-th value in $\Val(X)$ \\
$x^0,\ x^1$ (Boolean) & by convention $x^0 = \vn{false}$, $x^1 = \vn{true}$ \\
$\bx, \be, \ldots$ & a joint assignment of values to a set, e.g.\ $\bE = \be$ \\
$\Val(\bX)$ & all joint assignments to $\bX$: $\Val(X_1) \times \Val(X_2) \times \cdots$ \\
$\bx\langle\bY\rangle$ & assignment $\bx$ restricted to the variables in $\bY \subset \bX$ \\
$X_i^{(0)}, X_i^{(1)}, X_i^{(t)}$ & the template variable $X_i$ at time $t$ (temporal models, \cref{ch:dbn}) \\
$\params$ & the set of parameters of a model \\
$P(X = x^i)$ & a probability (value of the probability mass function) \\
$P(x^i),\ P(x^i, y^j)$ & shorthand for $P(X = x^i)$ and $P(X = x^i, Y = y^j)$ \\
$P(X, Y)$ & a probability distribution over the discrete variables $X, Y$ \\
$p(x)$ & a probability \emph{density} function for a continuous $X$ \\
\bottomrule
\end{tabular}
\caption{Notational conventions of the lecture. For named Boolean variables
the slides also write $\vn{flu}$ for $\vn{Flu} = \vn{true}$ and
$\neg\vn{flu}$ for $\vn{Flu} = \vn{false}$.}
\label{tab:notation}
\end{table}

% ----------------------------------------------------------
\section{Probability Distributions}
% ----------------------------------------------------------

\begin{defbox}[title={Probability distribution}]
Let $S$ be a set of mutually exclusive atomic events. A \emph{probability
distribution} over $S$ is a function $P: S \to \R$ with
\begin{enumerate}
  \item $P(\alpha) \geq 0$ for all $\alpha \in S$ \hfill (non-negativity)
  \item $P(\Omega) = 1$, where $\Omega = \bigcup_{\alpha_i \in S} \alpha_i$ is
    the disjunction of all possible events \hfill (normalisation)
  \item if $\alpha, \beta \in S$ and $\alpha \cap \beta = \emptyset$, then
    $P(\alpha \cup \beta) = P(\alpha) + P(\beta)$ \hfill (additivity)
\end{enumerate}
The \emph{probability} $P(\alpha)$ is the value the distribution assigns to the
event $\alpha$. Here $\cap$ and $\cup$ denote logical AND and OR between
events.
\end{defbox}

A tiny example with only two Boolean variables, $\vn{Flu}$ and $\vn{MP}$
(muscle pain), has four atomic events, and one possible distribution over them
is:
\begin{center}
\small
\begin{tabular}{@{}lcc@{}}
\toprule
$P(\vn{Flu}, \vn{MP})$ & $\vn{MP} = \vn{false}$ & $\vn{MP} = \vn{true}$ \\
\midrule
$\vn{Flu} = \vn{false}$ & 0.70 & 0.15 \\
$\vn{Flu} = \vn{true}$  & 0.05 & 0.10 \\
\bottomrule
\end{tabular}
\end{center}
So the probability that a randomly chosen person has the flu \emph{and} muscle
pain is $P(\vn{flu}, \vn{mp}) = 0.1$, and so on.

\begin{notebox}[title={Probability or distribution? Look at the argument}]
The same letter $P$ is used for two different kinds of objects:
\begin{itemize}
  \item if the argument is an \textbf{event} (values fixed), it is a
    \emph{probability} --- a single number: $P(\vn{flu}) = 0.15$;
  \item if the argument is a (set of) \textbf{random variable(s)}, it is a
    \emph{distribution} --- a whole table of numbers:
    $P(\vn{Flu}) = (0.85, 0.15)$.
\end{itemize}
Mixed cases follow the same rule. $P(\vn{Flu} \given \vn{mp})$ is one
distribution over $\vn{Flu}$, whereas $P(\vn{Flu} \given \vn{MP})$ is a
\emph{family} of distributions, one for every value of $\vn{MP}$.
\end{notebox}

\subsection{What does a probability mean?}

What exactly ``$P(\alpha) = 0.3$'' means is a deep philosophical question.
Intuitively, $P(\alpha)$ quantifies how confident we are that $\alpha$ is true
or will happen: $P(\alpha) = 1$ means we are certain it will, $P(\alpha) = 0$
that it will not. Everything in between needs an interpretation, and there are
two main camps:
\begin{itemize}
  \item \textbf{Frequentist:} $P(\alpha)$ is the fraction of times $\alpha$
    would occur if we repeated the experiment infinitely often. Fine for dice,
    awkward for ``the probability of rain \emph{this afternoon} is 0.7'' ---
    this afternoon happens only once.
  \item \textbf{Subjectivist (Bayesian):} $P(\alpha)$ is a degree of belief
    that $\alpha$ is true, or the degree to which the available evidence
    supports $\alpha$.
\end{itemize}
The course adopts the subjectivist view: $P(x)$ is the \emph{system's degree of
belief} that $x$ holds in the world. When we \emph{learn} probabilities from
data (\cref{ch:learning}), we will happily fall back on frequencies.

% ----------------------------------------------------------
\section{Joint and Marginal Distributions}
% ----------------------------------------------------------

\begin{defbox}[title={Full joint and marginal distributions}]
In a world described by the variables $\Xall$:
\begin{itemize}
  \item the \emph{full joint distribution} over $\Xall$ is the distribution
    over all atomic events, i.e.\ it assigns a probability to every
    combination of values of all variables;
  \item a \emph{marginal distribution} over a subset $\bX \subset \Xall$
    assigns probabilities to all value combinations of just the variables in
    $\bX$;
  \item the \emph{marginal distribution of a variable} $X \in \Xall$ is the
    special case $\bX = \{X\}$.
\end{itemize}
\end{defbox}

A few notational shortcuts: we write $P(X = x, Y = y)$ or just $P(x, y)$ for
the probability of the conjunction $(X = x) \cap (Y = y)$; $P(\bX) =
P(X_1, \ldots, X_k)$ for the joint distribution over a set; $P(\bX, \bY)$ for
the joint over $\bX \cup \bY$; and $P(\bX \given \bY)$ for the distribution
over $\bX$ conditioned on values of $\bY$ (next section).

The workhorse example of the slides is a diagnosis world with three Boolean
variables, $\Xall = \{\vn{Flu}, \vn{MP}, \vn{Fever}\}$, and this full joint
distribution:
\begin{equation}
\label{eq:flu-mp-fever}
\begin{array}{lcccc}
\toprule
 & \multicolumn{2}{c}{\vn{mp}} & \multicolumn{2}{c}{\neg\vn{mp}} \\
\cmidrule(lr){2-3} \cmidrule(lr){4-5}
P(\vn{Flu}, \vn{MP}, \vn{Fever}) & \vn{fever} & \neg\vn{fever} & \vn{fever} & \neg\vn{fever} \\
\midrule
\vn{flu}      & 0.15 & 0.10 & 0.10 & 0.03 \\
\neg\vn{flu}  & 0.03 & 0.04 & 0.05 & 0.50 \\
\bottomrule
\end{array}
\end{equation}
All eight entries sum to 1, as they must.

Marginals are obtained by \emph{summing out} everything you are not interested
in. To get $P(\vn{flu})$, add up all cells in the $\vn{flu}$ row --- whatever
the values of $\vn{MP}$ and $\vn{Fever}$:
\[
P(\vn{flu}) = \sum_{x, y} P(\vn{flu}, \vn{MP} = x, \vn{Fever} = y)
= 0.15 + 0.10 + 0.10 + 0.03 = 0.38 .
\]

\begin{exbox}[title={All three marginals of table \eqref{eq:flu-mp-fever}}]
\begin{align*}
P(\neg\vn{flu}) &= 0.03 + 0.04 + 0.05 + 0.50 = 0.62
  & P(\vn{Flu}) &= (0.62,\ 0.38) \\
P(\neg\vn{mp}) &= 0.10 + 0.03 + 0.05 + 0.50 = 0.68
  & P(\vn{MP}) &= (0.68,\ 0.32) \\
P(\neg\vn{fever}) &= 0.10 + 0.04 + 0.03 + 0.50 = 0.67
  & P(\vn{Fever}) &= (0.67,\ 0.33)
\end{align*}
Each distribution is written as $(P(x^0), P(x^1))$, i.e.\ \emph{false first}.
\end{exbox}

Why is summing allowed? The event $\vn{flu}$ is the disjunction ($\cup$) of all
atomic events with $\vn{Flu} = \vn{true}$, and atomic events are mutually
exclusive, so by additivity their probabilities simply add up. The same
argument works for \emph{any} event: \emph{the probability of an event
$\alpha$ is the sum over all atomic events (entries of the full joint) that
match $\alpha$}, i.e.\ that agree with $\alpha$ on the variables $\alpha$ fixes.
This one sentence is already the core of a general inference algorithm
(\cref{sec:enumeration}).

% ----------------------------------------------------------
\section{Conditional Probability, the Chain Rule and Bayes' Rule}
% ----------------------------------------------------------

Most questions we ask a model have the form ``I have observed $\beta$ --- what
is the probability of the (unobservable) thing $\alpha$ that I care about?''
For instance: given that my patient has a congestion, how likely is hayfever,
$P(\vn{Hayfever} = \vn{true} \given \vn{Congestion} = \vn{true})$? And if I
also learn that she has muscle pain, how does that change my belief? Or, for a
robot: given the current GPS and radar readings, what is the distribution over
its possible locations and velocities?

\begin{defbox}[title={Conditional probability}, unbreakable]
For two events $\alpha, \beta$ with $P(\beta) > 0$, the \emph{conditional
probability} of $\alpha$ given that $\beta$ is true is
\[
  P(\alpha \given \beta) = \frac{P(\alpha \cap \beta)}{P(\beta)} .
\]
\end{defbox}

The picture behind this formula (\cref{fig:conditioning}): learning that
$\beta$ is true rules out every world in which $\beta$ is false. What remains
is the region $\beta$, and inside it the worlds where $\alpha$ also holds. We
simply rescale so that the surviving region has total probability one.

\begin{figure}[htbp]
\centering
\begin{tikzpicture}[font=\small]
  % left panel: the whole space
  \def\betaE{(2.65,1.45) ellipse (1.45 and 0.95)}
  \def\alphaE{(1.45,1.45) ellipse (1.05 and 0.75)}
  \draw[accent, thick] (0,0) rectangle (4.6,2.9);
  \node[anchor=north west] at (0.05,2.88) {$\Omega$};
  \fill[figblue!12] \betaE;
  \begin{scope} \clip \betaE; \fill[figblue!45] \alphaE; \end{scope}
  \draw[figblue, thick] \betaE;
  \draw[figred, thick] \alphaE;
  \node[figred] at (0.75,1.45) {$\alpha$};
  \node[figblue] at (3.55,1.45) {$\beta$};
  \node[font=\scriptsize] at (1.95,1.45) {$\alpha\cap\beta$};
  % arrow
  \draw[-{Stealth[length=2.5mm]}, thick, accent] (4.9,1.45) -- (6.3,1.45)
    node[midway, above, align=center, font=\scriptsize] {observe $\beta$};
  % right panel: beta becomes the new universe
  \begin{scope}[shift={(6.7,-0.35)}, scale=1.25]
    \def\betaR{(1.75,1.45) ellipse (1.45 and 0.95)}
    \def\alphaR{(0.55,1.45) ellipse (1.05 and 0.75)}
    \fill[figblue!12] \betaR;
    \begin{scope} \clip \betaR; \fill[figblue!45] \alphaR; \end{scope}
    \draw[figblue, thick] \betaR;
    \node[font=\scriptsize] at (1.05,1.45) {$\alpha\cap\beta$};
    \node[figblue] at (2.65,1.45) {$\beta$};
  \end{scope}
  \node[align=center, font=\small] at (8.9,-0.25)
    {$P(\alpha \given \beta) = \dfrac{P(\alpha\cap\beta)}{P(\beta)}$};
\end{tikzpicture}
\caption{Conditioning as zooming in. Observing $\beta$ throws away all worlds
outside $\beta$; the region $\beta$ becomes the new universe and is rescaled
to total probability one. $P(\alpha \given \beta)$ is the share of $\beta$
that is also $\alpha$.}
\label{fig:conditioning}
\end{figure}

Two consequences are worth stating explicitly. For a fixed event $\beta$ and
some variables $\bX$, the conditional distribution $P(\bX \given \beta)$ has
all properties of a proper distribution --- in particular it sums to one over
all value combinations of $\bX$. And conditioning is an \emph{operation}: it
takes a distribution $P(\bX)$ and returns a new one, $P(\bX \given \beta)$.

In the two-variable table above, $P(\vn{flu}) = 0.05 + 0.10 = 0.15$, but
\[
P(\vn{flu} \given \vn{mp}) = \frac{P(\vn{flu}, \vn{mp})}{P(\vn{mp})}
= \frac{0.10}{0.15 + 0.10} = 0.4 .
\]
Learning about muscle pain raises the belief in flu from 15\,\% to 40\,\%.

\subsection{The chain rule}

Rearranging the definition of conditional probability gives the
\emph{product rule}, $P(\alpha \cap \beta) = P(\alpha \given \beta)\,P(\beta)$.
Because $\cap$ is symmetric, we may just as well write
$P(\alpha \cap \beta) = P(\alpha)\,P(\beta \given \alpha)$. Applying the
product rule over and over to $n$ events gives the chain rule.

\begin{thmbox}[title={Chain rule}, unbreakable]
\begin{align*}
P(\alpha_1, \ldots, \alpha_n)
&= P(\alpha_1 \given \alpha_2, \ldots, \alpha_n)\, P(\alpha_2, \ldots, \alpha_n) \\
&= P(\alpha_1 \given \alpha_2, \ldots, \alpha_n)\,
  P(\alpha_2 \given \alpha_3, \ldots, \alpha_n) \cdots
  P(\alpha_{n-1} \given \alpha_n)\, P(\alpha_n).
\end{align*}
Since conjunction is commutative and associative, any order works, e.g.
\[
P(\alpha_1, \ldots, \alpha_n) = P(\alpha_1)\,P(\alpha_2 \given \alpha_1) \cdots
P(\alpha_n \given \alpha_1, \ldots, \alpha_{n-1}).
\]
It also works partially,
\[
P(\alpha_1, \ldots, \alpha_n) = P(\alpha_1, \ldots, \alpha_k)\,
P(\alpha_{k+1}, \ldots, \alpha_n \given \alpha_1, \ldots, \alpha_k),
\]
and for whole distributions:
\[
P(X_1, \ldots, X_N) = P(X_1, \ldots, X_k)\,
P(X_{k+1}, \ldots, X_N \given X_1, \ldots, X_k).
\]
\end{thmbox}

The chain rule is \emph{always} true --- it makes no assumptions whatsoever. It
just says you can tell the story of a joint event one piece at a time, each new
piece conditioned on everything told so far. On its own it does not save
anything: the last factor still conditions on all the other variables. The
savings come later, when independence lets us shorten those conditioning sets
(\cref{sec:cond-indep} and \cref{ch:bn}).

\subsection{The law of total probability}

\begin{thmbox}[title={Law of total probability}]
For any event $\bX = \bx$ and any set of variables $\bY$:
\[
P(\bx) = \sum_{\by \in \Val(\bY)} P(\bx, \by)
       = \sum_{\by \in \Val(\bY)} P(\bx \given \by)\, P(\by).
\]
The same holds for whole distributions and with extra conditioning:
\[
P(\bX) = \sum_{\by} P(\bX \given \by)\, P(\by), \qquad
P(\bX \given \bZ) = \sum_{\by} P(\bX, \by \given \bZ)
                  = \sum_{\by} P(\bX \given \by, \bZ)\, P(\by \given \bZ).
\]
\end{thmbox}

The first form is exactly the ``summing out'' we used for marginals --- the law
of total probability is what justifies marginalisation and, below, inference
by enumeration.\footnote{The slides also list $P(\bx) = \sum_{\by}
P(\bx)\,P(\by \given \bx)$. That is true but trivial, because
$\sum_{\by} P(\by \given \bx) = 1$.} The second form is a \emph{case
analysis}: split the world into the cases $\bY = \by$, compute the probability
of $\bx$ in each case, and average, weighting each case by how likely it is.
With the two-variable table:
$P(\vn{mp}) = P(\vn{mp} \given \vn{flu})\,P(\vn{flu}) + P(\vn{mp} \given
\neg\vn{flu})\,P(\neg\vn{flu}) = \tfrac{0.10}{0.15} \cdot 0.15 +
\tfrac{0.15}{0.85} \cdot 0.85 = 0.25$.

\subsection{Bayes' rule}

Write the product rule both ways round, $P(\alpha \cap \beta) =
P(\alpha \given \beta)\,P(\beta) = P(\beta \given \alpha)\,P(\alpha)$, and
divide by $P(\beta)$:

\begin{thmbox}[title={Bayes' rule}, unbreakable]
\[
P(\alpha \given \beta) = \frac{P(\beta \given \alpha)\, P(\alpha)}{P(\beta)} .
\]
It computes a conditional probability from the ``inverse'' conditional
probability. In a diagnosis setting:
\[
\underbrace{P(\vn{Problem} \given \vn{Symptoms})}_{\text{posterior}}
= \frac{\overbrace{P(\vn{Symptoms} \given \vn{Problem})}^{\text{likelihood}}
        \;\overbrace{P(\vn{Problem})}^{\text{prior}}}
       {\underbrace{P(\vn{Symptoms})}_{\text{evidence}}}
\]
with
\[
P(\vn{Symptoms}) = \sum_{p \in \Val(\vn{Problem})} P(\vn{Symptoms} \given p)\, P(p).
\]
\end{thmbox}

Bayes' rule (after the Reverend Thomas Bayes, 1702--1761, whose essay was
published posthumously in 1763) is one of the most fundamental tools in
scientific data analysis and probabilistic reasoning. The reason is that
knowledge usually comes in the ``causal'' direction --- how often does flu
produce muscle pain? --- while questions come in the ``diagnostic'' direction:
this person has muscle pain, so do they have flu? The names of the four parts:

\begin{itemize}
  \item \textbf{Prior} $P(\vn{Problem})$: belief in the problem \emph{before}
    looking at the current case.
  \item \textbf{Posterior} $P(\vn{Problem} \given \vn{Symptoms})$: belief
    \emph{after} observing the symptoms.
  \item \textbf{Likelihood} $P(\vn{Symptoms} \given \vn{Problem})$: how likely
    the problem is to produce these symptoms.
  \item \textbf{Evidence} $P(\vn{Symptoms})$: how likely these symptoms are at
    all. By the law of total probability it is the sum above, so it can always
    be computed from priors and likelihoods.
\end{itemize}

\begin{exbox}[title={Does the person with muscle pain have the flu?}, unbreakable]
Suppose we know (or believe) that $P(\vn{mp} \given \vn{flu}) = 0.5$ (half of
all flu patients have muscle pain), $P(\vn{flu}) = 0.004$ (0.4\,\% of the
population has the flu at any time) and $P(\vn{mp}) = 0.1$ (10\,\% of people
have muscle pain). Then
\[
P(\vn{flu} \given \vn{mp}) = \frac{P(\vn{mp} \given \vn{flu})\,P(\vn{flu})}{P(\vn{mp})}
= \frac{0.5 \times 0.004}{0.1} = 0.02 .
\]
Only 2\,\%! Counting people makes this intuitive. Out of 100\,000 people, 400
have the flu, and 200 of them have muscle pain. But 10\,000 people have muscle
pain in total, so only $200 / 10\,000 = 2\,\%$ of those with muscle pain have
the flu. The tiny prior dominates.
\end{exbox}

\begin{notebox}[title={Do not flip the conditional}, unbreakable]
$P(\vn{mp} \given \vn{flu}) = 0.5$ but $P(\vn{flu} \given \vn{mp}) = 0.02$:
the two directions can be wildly different, and confusing them is the classic
mistake. Also note that a likelihood $P(\vn{Symptoms} \given \vn{Problem})$ is
a distribution over \emph{symptoms}, not over problems: summing it over all
problems does \emph{not} have to give one.
\end{notebox}

% ----------------------------------------------------------
\section{Probabilistic Inference: Queries and Enumeration}
\label{sec:enumeration}
% ----------------------------------------------------------

Here is the scenario for the rest of the chapter. The full joint distribution
over $\Xall$ is the system's \emph{knowledge base}: its model of the world,
assigning a probability to every possible state. The system's job is to
compute, from this knowledge base and the observed values of some variables,
the probability of any hypothesis we are interested in. Deriving the truth or
probability of a hypothesis from a knowledge base and observations is called
\emph{inference} or \emph{reasoning}. For us it almost always means computing
(conditional) distributions over the variables of interest.

\begin{defbox}[title={Probabilistic query}]
A \emph{probabilistic query} asks for the conditional distribution
\[
  P(\bX \given \bE = \be), \quad\text{or } P(\bX \given \be) \text{ for short},
\]
for a set of \emph{query variables} $\bX \subseteq \Xall$, a set of
\emph{evidence variables} $\bE \subseteq \Xall$ and their observed values $\be$.
In words: given that we have observed $\be$, how probable is each value
combination of $\bX$? All remaining variables are \emph{hidden}.
\end{defbox}

With table \eqref{eq:flu-mp-fever} we could ask for an event probability,
$P(\vn{fever} \wedge \neg\vn{mp})$ (someone has a fever but no muscle pain); a
marginal distribution, $P(\vn{Flu})$; or a conditional distribution,
$P(\vn{Flu} \given \vn{fever}, \neg\vn{mp})$. All of them are answered the
same way.

\subsection{Inference by enumeration}

Event probabilities are just sums of matching cells:
\begin{align*}
P(\neg\vn{fever}) &= \sum_{x, y} P(\neg\vn{fever}, \vn{MP} = x, \vn{Flu} = y)
  = 0.10 + 0.03 + 0.04 + 0.50 = 0.67, \\
P(\vn{fever} \wedge \neg\vn{mp}) &= \sum_{x \in \Val(\vn{Flu})}
  P(\vn{fever}, \neg\vn{mp}, \vn{Flu} = x) = 0.10 + 0.05 = 0.15 .
\end{align*}
Conditional probabilities are a ratio of two such sums:
\begin{align*}
P(\vn{flu} \given \neg\vn{fever})
&= \frac{P(\vn{flu} \wedge \neg\vn{fever})}{P(\neg\vn{fever})}
= \frac{0.10 + 0.03}{0.10 + 0.03 + 0.04 + 0.50}
= \frac{0.13}{0.67} \approx 0.194, \\
P(\neg\vn{flu} \given \neg\vn{fever})
&= \frac{P(\neg\vn{flu} \wedge \neg\vn{fever})}{P(\neg\vn{fever})}
= \frac{0.04 + 0.50}{0.67} = \frac{0.54}{0.67} \approx 0.806 .
\end{align*}

Look at the denominators: both are $P(\neg\vn{fever}) = 0.67$, because the
evidence probability does not depend on the value of the query variable. In
general, when computing $P(X \given y)$ we would compute the same $P(y)$
$|\Val(X)|$ times. There is a smarter way: $P(y)$ is just the constant that
makes the result sum to one, a \emph{normalisation constant} $Z$, and we can
get it for free afterwards by adding up the unnormalised numbers.

\begin{thmbox}[title={Inference by enumeration}, unbreakable]
To compute $P(\bX \given \be)$, let $\bZ = \Xall \setminus (\bX \cup \bE)$ be the
hidden variables.
\begin{enumerate}
  \item \textbf{Sum out the hidden variables.} For every $\bx \in \Val(\bX)$,
    \[ P(\bx, \be) = \sum_{\bz \in \Val(\bZ)} P(\bx, \be, \bz). \]
  \item \textbf{Compute the normalisation constant} by adding up the results
    of step 1 (law of total probability):
    \[ Z = P(\be) = \sum_{\bx} P(\bx, \be). \]
  \item \textbf{Renormalise:}
    \[ P(\bX \given \be) = \frac{1}{Z}\, P(\bX, \be). \]
\end{enumerate}
\end{thmbox}

\begin{exbox}[title={Two queries, step by step}]
\textbf{Query} $P(\vn{Flu} \given \neg\vn{fever})$; the hidden variable is
$\vn{MP}$. Vectors list $(\vn{flu}, \neg\vn{flu})$ top to bottom:
\[
P(\vn{Flu}, \neg\vn{fever}) = \sum_{x \in \{\vn{mp}, \neg\vn{mp}\}}
P(\vn{Flu}, \neg\vn{fever}, x)
= \begin{pmatrix} 0.10 + 0.03 \\ 0.04 + 0.50 \end{pmatrix}
= \begin{pmatrix} 0.13 \\ 0.54 \end{pmatrix},
\]
\[
Z = 0.13 + 0.54 = 0.67,
\qquad P(\vn{Flu} \given \neg\vn{fever}) = \frac{1}{Z}
\begin{pmatrix} 0.13 \\ 0.54 \end{pmatrix} \approx
\begin{pmatrix} 0.194 \\ 0.806 \end{pmatrix}.
\]
\textbf{Query} $P(\vn{Flu} \given \vn{fever}, \neg\vn{mp})$; nothing is hidden,
so step 1 just reads off two cells:
\[
P(\vn{Flu}, \vn{fever}, \neg\vn{mp}) = \begin{pmatrix} 0.10 \\ 0.05 \end{pmatrix},
\quad Z = 0.15,
\quad P(\vn{Flu} \given \vn{fever}, \neg\vn{mp}) \approx
\begin{pmatrix} 0.667 \\ 0.333 \end{pmatrix}.
\]
A fever without muscle pain makes flu twice as likely as no flu.
\end{exbox}

\begin{notebox}[title={Always label the entries of a distribution}, unbreakable]
The slides write the same result once as $P(\vn{Flu} \given \neg\vn{fever}) =
\{0.806, 0.194\}$ (order $\neg\vn{flu}, \vn{flu}$) and once as $\{0.194,
0.806\}$ (order $\vn{flu}, \neg\vn{flu}$). Neither is wrong, but an unlabeled
vector is ambiguous. On the exam, either label the entries or stick to the
convention $x^0$ (false) first.
\end{notebox}

Inference by enumeration is completely general: any query, any evidence, one
algorithm. Its problem is cost. With $n$ variables of $d$ values each, the
joint table has $d^n$ entries, and the sum in step 1 may have to touch almost
all of them. Exponential space \emph{and} exponential time --- this is the
first concrete appearance of the complexity problems from
\cref{sec:diagnosis-world}, and the reason the next sections are about
independence.

\subsection{Warm-up: everything you can ask of a small table}

The slides close this part with a self-test (motivated by a tutorial by Ali
Taylan Cemgil). It introduces a few quantities --- maximum, mode and
max-marginal --- that will come back when we look for the \emph{most likely}
explanation of some evidence (\cref{ch:hmm}).

\begin{exbox}[title={Self-test with $X_1, X_2 \in \{1, 2\}$}]
\begin{center}
\small
\begin{tabular}{@{}lcc@{}}
\toprule
$P(X_1, X_2)$ & $X_2 = 1$ & $X_2 = 2$ \\
\midrule
$X_1 = 1$ & 0.3 & 0.2 \\
$X_1 = 2$ & 0.1 & 0.4 \\
\bottomrule
\end{tabular}
\end{center}
\begin{itemize}
  \item \textbf{Marginals:} $P(X_1) = (0.5,\ 0.5)$ (row sums);
    $P(X_2) = (0.4,\ 0.6)$ (column sums).
  \item \textbf{Conditionals:} divide each column by its sum for
    $P(X_1 \given X_2)$, each row by its sum for $P(X_2 \given X_1)$:
    \[
    \begin{array}{l|cc}
    P(X_1 \given X_2) & X_2 = 1 & X_2 = 2 \\ \hline
    X_1 = 1 & 0.75 & 0.33 \\
    X_1 = 2 & 0.25 & 0.67
    \end{array}
    \qquad
    \begin{array}{l|cc}
    P(X_2 \given X_1) & X_2 = 1 & X_2 = 2 \\ \hline
    X_1 = 1 & 0.6 & 0.4 \\
    X_1 = 2 & 0.2 & 0.8
    \end{array}
    \]
  \item \textbf{Joint vs.\ posterior:} $P(X_1, X_2 = 2) = (0.2,\ 0.4)$ is a
    slice of the table and does \emph{not} sum to one; $P(X_1 \given X_2 = 2)
    = (0.33,\ 0.67)$ is that slice, normalised.
  \item \textbf{Evidence:} the normaliser of that slice, $P(X_2 = 2) =
    \sum_{x_1} P(x_1, X_2 = 2) = 0.6$.
  \item \textbf{Normalisation constant} of the whole table:
    $P(\Omega) = \sum_{x_1}\sum_{x_2} P(x_1, x_2) = 1$.
  \item \textbf{Max:} $\max_{x_1} P(x_1 \given X_2 = 1) = 0.75$ --- the
    \emph{value} of the best probability.
  \item \textbf{Mode:} $x_1^* = \argmax_{x_1} P(x_1 \given X_2 = 1) = 1$ ---
    \emph{which} value is best.
  \item \textbf{Max-marginal:} $\max_{x_1} P(x_1, X_2) = (0.3,\ 0.4)$, the
    ``skyline'' of the table: for every $x_2$, the largest entry in its column.
    Unlike a marginal it maximises instead of summing.
  \item \textbf{Independent?} No: $P(X_1 = 1 \given X_2 = 1) = 0.75 \neq 0.5 =
    P(X_1 = 1)$. Equivalently, $P(1, 1) = 0.3 \neq P(X_1 = 1)\,P(X_2 = 1) =
    0.2$.
\end{itemize}
\end{exbox}

% ----------------------------------------------------------
\section{Independence}
\label{sec:indep}
% ----------------------------------------------------------

Normally, learning that one event happened changes our belief in another:
$P(\alpha \given \beta) \neq P(\alpha)$. Sometimes it does not. Toss two coins:
whether the second one shows heads is not affected at all by what the first
one showed. That is independence.

\begin{defbox}[title={Independence of events}]
Two events $\alpha$ and $\beta$ are \emph{independent} in $P$, written
$(\alpha \indep \beta)$, if (equivalently)
\[
P(\alpha \given \beta) = P(\alpha), \qquad
P(\beta \given \alpha) = P(\beta), \qquad
P(\alpha \cap \beta) = P(\alpha)\,P(\beta).
\]
Independence is symmetric: $(\alpha \indep \beta) \Leftrightarrow (\beta \indep
\alpha)$.
\end{defbox}

The slides leave the equivalence as an exercise; it is two lines (assuming
$P(\alpha), P(\beta) > 0$). If $P(\alpha \given \beta) = P(\alpha)$, the
product rule gives $P(\alpha \cap \beta) = P(\alpha \given \beta)\,P(\beta) =
P(\alpha)\,P(\beta)$. If $P(\alpha \cap \beta) = P(\alpha)\,P(\beta)$, then
$P(\beta \given \alpha) = P(\alpha \cap \beta)/P(\alpha) = P(\beta)$, and in
the same way $P(\alpha \given \beta) = P(\alpha)$. The product form is
symmetric in $\alpha$ and $\beta$, which proves the symmetry too.

\begin{defbox}[title={Independence of random variables (and sets of them)}]
Two sets of random variables $\bX$ and $\bY$ are \emph{independent}, written
$(\bX \indep \bY)$, if for \emph{all} value combinations $\bx \in \Val(\bX)$,
$\by \in \Val(\bY)$:
\[
P(\bx \given \by) = P(\bx) \quad\text{and}\quad P(\by \given \bx) = P(\by),
\]
in short $P(\bX \given \bY) = P(\bX)$ and $P(\bY \given \bX) = P(\bY)$.
Consequence: $P(\bX, \bY) = P(\bX)\,P(\bY)$. Single variables are the special
case $\bX = \{X\}$, $\bY = \{Y\}$.
\end{defbox}

Why does this matter so much? Because then every entry of the joint can be
rebuilt from two smaller tables, $P(\bx, \by) = P(\bx)\,P(\by)$, so we only
need to store distributions over $\bX$ and over $\bY$ separately instead of
over $\bX \cup \bY$. Since table size grows exponentially with the number of
variables, that is a big saving in space.

The extreme case: $N$ independent coins $X_1, \ldots, X_N$. Then
\[
P(X_1, \ldots, X_N) = P(X_1)\,P(X_2) \cdots P(X_N).
\]
The joint has $2^N$
entries, but each $P(X_i)$ has only two (and really only one number, since
$P(\text{tails}) = 1 - P(\text{heads})$). Storing $2N$ --- or $N$ ---
numbers instead of $2^N$ is an exponential reduction from $\bigO{2^N}$ to
$\bigO{N}$.

\begin{notebox}[title={Checking independence: one cell can only disprove it}, unbreakable]
Variable independence must hold for \emph{every} combination of values. A
single cell where $P(x, y) \neq P(x)\,P(y)$ proves dependence; a single cell
where the product happens to match proves nothing. In table
\eqref{eq:flu-mp-fever}, for example, $P(\vn{flu}, \vn{mp}) = 0.25$ but
$P(\vn{flu})\,P(\vn{mp}) = 0.38 \times 0.32 \approx 0.12$, so $\vn{Flu}$ and
$\vn{MP}$ are dependent --- as you would expect.
\end{notebox}

% ----------------------------------------------------------
\section{Conditional Independence}
\label{sec:cond-indep}
% ----------------------------------------------------------

Full independence is rare in interesting worlds; things \emph{are} connected.
What is common is a weaker, and much more useful, property. Take
\begin{center}
$\alpha$ = ``patient Smith has a fever'', \quad
$\beta$ = ``patient Smith has muscle pain'', \quad
$\gamma$ = ``patient Smith has the flu''.
\end{center}
Fever and muscle pain are dependent: they often go together, so learning about
muscle pain raises our expectation of fever, $P(\alpha \given \beta) >
P(\alpha)$. But suppose we already \emph{know} that Smith has the flu. Then
learning about her muscle pain tells us nothing new about the fever --- the
fever is already ``explained'' by the flu: $P(\alpha \given \beta \cap \gamma)
= P(\alpha \given \gamma)$.

\begin{defbox}[title={Conditional independence}]
Two events $\alpha, \beta$ are \emph{conditionally independent given} $\gamma$,
written $(\alpha \indep \beta \given \gamma)$, if (equivalently)
\[
P(\alpha \given \beta \cap \gamma) = P(\alpha \given \gamma), \qquad
P(\beta \given \alpha \cap \gamma) = P(\beta \given \gamma), \qquad
P(\alpha \cap \beta \given \gamma) = P(\alpha \given \gamma)\,P(\beta \given \gamma).
\]
For sets of variables, $(\bX \indep \bY \given \bZ)$ if for all $\bx, \by, \bz$:
$P(\bx \given \by, \bz) = P(\bx \given \bz)$ and $P(\by \given \bx, \bz) =
P(\by \given \bz)$, in short $P(\bX \given \bY, \bZ) = P(\bX \given \bZ)$.
Consequence: $P(\bX, \bY \given \bZ) = P(\bX \given \bZ)\,P(\bY \given \bZ)$.
Conditional independence is symmetric as well.
\end{defbox}

The equivalence proof is the same as for plain independence, with
``$\given \gamma$'' carried along everywhere. Now comes the step the entire
course is built on:

\begin{thmbox}[title={Factorising with conditional independence}]
If $(\bX \indep \bY \given \bZ)$, the joint distribution splits into smaller
pieces:
\begin{align*}
P(\bX, \bY, \bZ) &= P(\bX \given \bY, \bZ)\, P(\bY, \bZ)
  && \text{chain rule} \\
&= P(\bX \given \bY, \bZ)\, P(\bY \given \bZ)\, P(\bZ)
  && \text{chain rule} \\
&= P(\bX \given \bZ)\, P(\bY \given \bZ)\, P(\bZ)
  && \text{conditional independence.}
\end{align*}
\end{thmbox}

Make sure you understand these three lines --- they are the most fundamental
idea in graphical models. The chain rule always works but leaves long
conditioning sets; conditional independence lets us \emph{delete} variables
from them. A Bayesian network (\cref{ch:bn}) is nothing but this trick applied
systematically to every variable.

\begin{exbox}[title={Fever and muscle pain, with numbers}, unbreakable]
Build a joint over $\vn{Flu}, \vn{Fever}, \vn{MP}$ from the factorisation
$P(\vn{Flu})\,P(\vn{Fever} \given \vn{Flu})\,P(\vn{MP} \given \vn{Flu})$ with
\begin{gather*}
P(\vn{flu}) = 0.1, \\
P(\vn{fever} \given \vn{flu}) = 0.8,\quad P(\vn{fever} \given \neg\vn{flu}) = 0.1,\qquad
P(\vn{mp} \given \vn{flu}) = 0.7,\quad P(\vn{mp} \given \neg\vn{flu}) = 0.2 .
\end{gather*}
\textbf{Marginally dependent.} By total probability $P(\vn{fever}) = 0.1 \cdot
0.8 + 0.9 \cdot 0.1 = 0.17$, $P(\vn{mp}) = 0.1 \cdot 0.7 + 0.9 \cdot 0.2 =
0.25$, and $P(\vn{fever}, \vn{mp}) = 0.1 \cdot 0.8 \cdot 0.7 + 0.9 \cdot 0.1
\cdot 0.2 = 0.074$. Hence
\[
P(\vn{fever} \given \vn{mp}) = 0.074 / 0.25 = 0.296 \;>\; 0.17 = P(\vn{fever}).
\]
Muscle pain is evidence for flu, and flu is evidence for fever.\par
\textbf{Conditionally independent.} Once flu is known, muscle pain adds
nothing:
\[
P(\vn{fever} \given \vn{mp}, \vn{flu}) = 0.8 = P(\vn{fever} \given \vn{flu}).
\]
\textbf{Saving.} The factorised form needs $1 + 2 + 2 = 5$ numbers instead of
the $2^3 - 1 = 7$ of a general joint. Tiny here, enormous for big models.
\end{exbox}

\begin{intbox}[unbreakable]
Once you know the cause, its effects stop ``talking'' to each other. All of
the correlation between fever and muscle pain was routed \emph{through} the
flu; fixing the flu cuts that route. This is why causal stories are such a
good guide to independencies: effects of a common cause are independent given
that cause.
\end{intbox}

\begin{notebox}[title={Independence and conditional independence do not imply each other}]
\begin{itemize}
  \item \textbf{Dependent, but independent given $Z$:} fever and muscle pain
    above (common cause flu).
  \item \textbf{Independent, but dependent given $Z$:} toss two fair coins
    independently and let $Z$ = ``both coins show the same face''. Given $Z$,
    the first coin \emph{determines} the second, so they are dependent. We will
    meet this pattern as ``explaining away'' in \cref{ch:bn}.
\end{itemize}
\end{notebox}

% ----------------------------------------------------------
\section{Continuous Random Variables}
\label{sec:continuous}
% ----------------------------------------------------------

\subsection{Probabilities versus densities}

So far every variable had a finite domain. Real applications, however, are
full of continuous quantities: range sensors, temperature and pressure, GPS,
camera pixels. And here a problem appears: each individual value has
probability zero. Take $\Val(X) = [0, 1]$ and give every number the same
probability. There are infinitely many numbers, so if that probability were
any positive value, the total would be infinite. Hence $P(X = x) = 0$ for all
$x$. (The same happens for non-uniform distributions.) We need a different
object.

\begin{defbox}[title={Probability density and cumulative distribution}, unbreakable]
A function $p: \R \to \R$ is a \emph{probability density function} (PDF) for a
continuous variable $X$ if $p(x) \geq 0$ for all $x \in \Val(X)$, $p(x) = 0$
for all $x \notin \Val(X)$, and $\int_{-\infty}^{+\infty} p(x) \dd x = 1$.
The \emph{cumulative distribution function} (CDF) is
\[
P(X \leq a) = \int_{-\infty}^{a} p(x) \dd x,
\]
so that
\[
P(a \leq X \leq b) = P(X \leq b) - P(X \leq a) = \int_{a}^{b} p(x) \dd x .
\]
\end{defbox}

A PDF is \emph{not} a probability distribution, and $p(x)$ is \emph{not} a
probability --- it can even be larger than one. Probabilities live on
intervals and are computed by integrating the density. The best mental
picture: $p(x)$ is probability \emph{per unit length}, the rate at which the
CDF grows at $x$. For a small interval of width $\delta$,
\[
P(x \leq X \leq x + \delta) \approx p(x)\,\delta .
\]

\subsection{The uniform distribution}

$X$ is \emph{uniformly distributed} on $[a, b]$, written $X \sim
\mathit{Unif}[a, b]$, if
\[
p(x) = \begin{cases} \dfrac{1}{b - a} & a \leq x \leq b, \\[4pt]
                     0 & \text{otherwise.} \end{cases}
\]
\Cref{fig:uniform} shows the PDF and the CDF. Note what happens if the
interval is short: on $[0, 0.5]$ the density is $2$ everywhere. That is
perfectly legal --- only the \emph{area} under a PDF has to be one.

\begin{figure}[htbp]
\centering
\begin{tikzpicture}
\begin{axis}[
  name=pdf, width=6.4cm, height=4.2cm, axis lines=left,
  xmin=0, xmax=4.2, ymin=0, ymax=0.75,
  xtick={1,3}, xticklabels={$a$,$b$},
  ytick={0.5}, yticklabels={$\frac{1}{b-a}$},
  xlabel={$x$}, ylabel={$p(x)$}, title={\small PDF},
  ]
  \addplot[figblue, very thick] coordinates {(0,0) (1,0)};
  \addplot[figblue, very thick] coordinates {(1,0.5) (3,0.5)};
  \addplot[figblue, very thick] coordinates {(3,0) (4.2,0)};
  \draw[dashed, gray] (axis cs:1,0) -- (axis cs:1,0.5);
  \draw[dashed, gray] (axis cs:3,0) -- (axis cs:3,0.5);
\end{axis}
\begin{axis}[
  at={(pdf.east)}, anchor=west, xshift=1.6cm,
  width=6.4cm, height=4.2cm, axis lines=left,
  xmin=0, xmax=4.2, ymin=0, ymax=1.2,
  xtick={1,3}, xticklabels={$a$,$b$}, ytick={1},
  xlabel={$x$}, ylabel={$P(X \leq x)$}, title={\small CDF},
  ]
  \addplot[figblue, very thick] coordinates {(0,0) (1,0) (3,1) (4.2,1)};
  \draw[dashed, gray] (axis cs:3,0) -- (axis cs:3,1);
\end{axis}
\end{tikzpicture}
\caption{The uniform distribution on $[a, b]$. The CDF grows linearly from 0 to
1 across the interval; its slope is the density.}
\label{fig:uniform}
\end{figure}

\subsection{The normal (Gaussian) distribution}

\begin{defbox}[title={Normal distribution}]
$X$ has a \emph{normal (Gaussian) distribution} with mean $\mu$ and variance
$\sigma^2$, written $X \sim \Normal(\mu, \sigma^2)$, if
\[
p(x) = \frac{1}{\sqrt{2\pi}\,\sigma}\; e^{-\frac{(x - \mu)^2}{2\sigma^2}} .
\]
\end{defbox}

The mean $\mu$ says where the bell is centred; the standard deviation $\sigma$
says how wide it is (\cref{fig:gaussians}). A handy rule of thumb: about
68\,\% of the probability mass lies within $\mu \pm \sigma$, about 95\,\%
within $\mu \pm 2\sigma$. Gaussians are everywhere in this course's temporal
part, because sums of many small random disturbances tend to look Gaussian and
because they are mathematically very convenient (\cref{ch:kalman}).

\begin{figure}[htbp]
\centering
\begin{tikzpicture}
\begin{axis}[
  width=11cm, height=5cm, axis lines=left,
  xmin=-10, xmax=10, ymin=0, ymax=0.45,
  xlabel={$x$}, ylabel={$p(x)$},
  every axis plot/.append style={very thick, samples=200, smooth},
  ytick={0,0.1,0.2,0.3,0.4}]
  \addplot[figblue, domain=-10:10]  {exp(-(x)^2/2)/sqrt(2*pi)};
  \addplot[figred, domain=-10:10]   {exp(-(x-5)^2/(2*4))/(sqrt(2*pi)*2)};
  \addplot[figgreen, domain=-10:10] {exp(-(x)^2/(2*16))/(sqrt(2*pi)*4)};
  \node[figblue, anchor=west]  at (axis cs:0.9,0.36)  {$\Normal(0, 1)$};
  \node[figred, anchor=south]  at (axis cs:5,0.205)   {$\Normal(5, 2^2)$};
  \node[figgreen, anchor=south] at (axis cs:-6.5,0.07) {$\Normal(0, 4^2)$};
\end{axis}
\end{tikzpicture}
\caption{Three Gaussian densities. Wider bells are lower, because the area
under each curve is always one.}
\label{fig:gaussians}
\end{figure}

\subsection{The multivariate Gaussian}

\begin{defbox}[title={Multivariate normal distribution}]
$N$ variables $\bX = \{X_1, \ldots, X_N\}$ have a \emph{joint (multivariate)
normal distribution} with mean vector $\bm{\mu}$ and covariance matrix
$\Sigma$, written $\bX \sim \Normal(\bm{\mu}, \Sigma)$, if
\[
p(\bx) = \frac{1}{(2\pi)^{N/2}\,\abs{\Sigma}^{1/2}}\;
\exp\!\Big[-\tfrac{1}{2}(\bx - \bm{\mu})\T\, \Sigma^{-1}\, (\bx - \bm{\mu})\Big],
\]
where $\bm{\mu} = (\mu_{X_1}, \ldots, \mu_{X_N})$, $\Sigma$ is the $N \times N$
covariance matrix with $\Sigma_{ij} = \sigma_{ij}$ the covariance between
$X_i$ and $X_j$, and $\abs{\Sigma}$ is its determinant.
\end{defbox}

In two dimensions the density is a ``hill'' whose contour lines are ellipses
(\cref{fig:gauss2d}). If $X$ and $Y$ are uncorrelated ($\sigma_{XY} = 0$), the
ellipses are aligned with the axes. If they are positively correlated
($\sigma_{XY} > 0$), the ellipses tilt: large $x$ tends to come with large
$y$. The slides call $\sigma_{ij}$ the ``pairwise correlation''; strictly it
is the \emph{covariance} --- the correlation coefficient is the covariance
divided by $\sigma_i \sigma_j$ --- but both are zero or non-zero together.

\begin{figure}[htbp]
\centering
\begin{tikzpicture}[font=\small]
  \foreach \dx/\rot/\ra/\rb/\lab in {0/0/1.15/0.85/{$\sigma_{XY} = 0$},
                                     6/35/1.45/0.6/{$\sigma_{XY} > 0$}} {
    \begin{scope}[xshift=\dx cm]
      \draw[-{Stealth[length=2mm]}, gray] (-2,0) -- (2.1,0) node[right] {$x$};
      \draw[-{Stealth[length=2mm]}, gray] (0,-1.7) -- (0,1.8) node[above] {$y$};
      \foreach \s/\op in {0.45/70, 0.85/45, 1.25/25} {
        \draw[figblue, thick, rotate=\rot, opacity=1]
          (0,0) ellipse ({\s*\ra} and {\s*\rb});
      }
      \fill[figred] (0,0) circle (1.6pt);
      \node at (0,-2.2) {\lab};
    \end{scope}
  }
\end{tikzpicture}
\caption{Contour lines of two 2-D Gaussians centred at $\bm{\mu} = \zeros$.
Left: uncorrelated variables, axis-aligned ellipses. Right: positively
correlated variables, tilted ellipses.}
\label{fig:gauss2d}
\end{figure}

\begin{notebox}[title={Uncorrelated vs.\ independent}, unbreakable]
For a \emph{joint Gaussian}, uncorrelated means independent: if $\Sigma$ is
diagonal, the density factorises into a product of one-dimensional Gaussians.
For general distributions this is false --- zero correlation only rules out a
\emph{linear} relationship.
\end{notebox}

% ----------------------------------------------------------
\section{Chapter Checklist}
% ----------------------------------------------------------

After this chapter you should be able to:
\begin{itemize}
  \item read the lecture's notation, and tell a probability from a
    distribution by looking at the argument of $P(\cdot)$;
  \item compute marginals, conditionals and event probabilities from a joint
    table, and use the chain rule, the law of total probability and Bayes'
    rule;
  \item state the probabilistic query $P(\bX \given \be)$ and answer it by
    enumeration, including the normalisation trick;
  \item define (conditional) independence for events and for sets of
    variables, check it in a table, and explain why it makes representations
    smaller;
  \item explain the difference between a density and a probability, and work
    with uniform and (multivariate) Gaussian densities.
\end{itemize}


% ==========================================================
%  CHAPTER 3 — BAYESIAN NETWORKS: REPRESENTATION   (part 3, VO 3)
% ==========================================================
\chapter[Bayesian Networks: Representation and Semantics]{Bayesian Networks:\\ Representation and Semantics}
\label{ch:bn}

\Cref{ch:prob} ended with a promise and a problem. The promise: any query can
be answered from the full joint distribution. The problem: the full joint is
exponentially large. We also saw the way out --- conditional independence lets
us write the joint as a product of small pieces. A \emph{Bayesian network}
(BN) is the data structure that does this systematically: a directed acyclic
graph that tells us \emph{which} small pieces we need, plus one small table
per variable that holds the numbers.

This chapter introduces BNs through one running example, the \emph{student
network}. We first play with it intuitively (which kinds of reasoning does it
support?), then make its meaning precise in two equivalent ways --- as a set of
independence statements and as a factorised joint distribution --- and finally
look at how compact it is, how to build one, and how to handle continuous
variables.

% ----------------------------------------------------------
\section{A Running Example: The Student Network}
% ----------------------------------------------------------

\subsection{The story}

The example comes from Koller \& Friedman, with a small twist by the lecturer.
A company wants to hire a college graduate. It is looking for genuinely
interested and motivated people, but motivation cannot be tested directly.
Instead, applicants hand in their SAT score (a standardised US university
admission test) and their grade in a relevant university course. The company
wants to use these two pieces of evidence to estimate whether an applicant is
highly motivated. Our job: turn this into a quantitative model that algorithms
can compute with.

\paragraph{Step 1: variables} Three concepts matter for the decision, each
modelled as a discrete random variable with a deliberately coarse domain:
\[
\begin{array}{lll}
\text{Interest/Motivation } (I) & \Val(I) = \{\vn{lo}, \vn{hi}\} & = \{i^0, i^1\} \\
\text{SAT score } (S)           & \Val(S) = \{\vn{lo}, \vn{hi}\} & = \{s^0, s^1\} \\
\text{Grade } (G)               & \Val(G) = \{A, B, C\}          & = \{g^1, g^2, g^3\}
\end{array}
\]
Careful with the grade: $g^1$ is the \emph{best} grade (A) and $g^3$ the worst
(C).

\paragraph{Step 2: dependency structure} The SAT score depends (among other
things) on motivation, and so does the course grade. The story then grows a
little. The grade also depends on how \emph{difficult} the course was. And the
student asks the professor for a recommendation \emph{letter}, which the
professor writes based on the course grade only --- there are too many
students to remember anyone personally. Two more variables:
\[
\begin{array}{lll}
\text{Difficulty } (D) & \Val(D) = \{\vn{easy}, \vn{hard}\}  & = \{d^0, d^1\} \\
\text{Letter } (L)     & \Val(L) = \{\vn{weak}, \vn{strong}\} & = \{l^0, l^1\}
\end{array}
\]
Drawing an arrow from each cause to its effect gives the graph in
\cref{fig:student}: $D \to G \leftarrow I$, $I \to S$, $G \to L$.

\paragraph{Step 3: uncertainty} None of these relationships is deterministic.
Motivation influences the SAT score but is not the only factor (and some SAT
tests are easier than others). Motivation and difficulty do not uniquely
determine the grade. An absent-minded professor might write a strong letter
despite a poor grade. We model all of this with probabilities: every variable
gets a \emph{conditional probability distribution} (CPD) given its parents.

\begin{figure}[htbp]
\centering
\begin{tikzpicture}
  % network
  \node[bn] (D) at (-1.9, 0)   {Difficulty};
  \node[bn] (I) at ( 1.9, 0)   {Interest};
  \node[bn] (G) at (-0.7,-1.9) {Grade};
  \node[bn] (S) at ( 3.1,-1.9) {SAT};
  \node[bn] (L) at (-0.7,-3.8) {Letter};
  \draw[edge] (D) -- (G);
  \draw[edge] (I) -- (G);
  \draw[edge] (I) -- (S);
  \draw[edge] (G) -- (L);
  % CPDs
  \node[cpd, anchor=south east,
        label={[font=\scriptsize\bfseries]above:$P(D)$}] (tD) at (-3.2,0.3)
    {\begin{tabular}{cc}\toprule $d^0$ & $d^1$\\\midrule 0.6 & 0.4\\\bottomrule\end{tabular}};
  \node[cpd, anchor=south west,
        label={[font=\scriptsize\bfseries]above:$P(I)$}] (tI) at (3.2,0.3)
    {\begin{tabular}{cc}\toprule $i^0$ & $i^1$\\\midrule 0.7 & 0.3\\\bottomrule\end{tabular}};
  \node[cpd, anchor=east,
        label={[font=\scriptsize\bfseries]above:$P(G \given I, D)$}] (tG) at (-2.3,-2.3)
    {\begin{tabular}{lccc}\toprule
       & $g^1$ & $g^2$ & $g^3$\\\midrule
     $i^0, d^0$ & 0.3  & 0.4  & 0.3 \\
     $i^0, d^1$ & 0.05 & 0.25 & 0.7 \\
     $i^1, d^0$ & 0.9  & 0.08 & 0.02 \\
     $i^1, d^1$ & 0.5  & 0.3  & 0.2 \\\bottomrule\end{tabular}};
  \node[cpd, anchor=west,
        label={[font=\scriptsize\bfseries]above:$P(S \given I)$}] (tS) at (4.3,-2.0)
    {\begin{tabular}{lcc}\toprule
       & $s^0$ & $s^1$\\\midrule
     $i^0$ & 0.95 & 0.05 \\
     $i^1$ & 0.2  & 0.8 \\\bottomrule\end{tabular}};
  \node[cpd, anchor=west,
        label={[font=\scriptsize\bfseries]above:$P(L \given G)$}] (tL) at (0.6,-4.3)
    {\begin{tabular}{lcc}\toprule
       & $l^0$ & $l^1$\\\midrule
     $g^1$ & 0.1  & 0.9 \\
     $g^2$ & 0.4  & 0.6 \\
     $g^3$ & 0.99 & 0.01 \\\bottomrule\end{tabular}};
  \draw[cpdlink] (tD) -- (D);
  \draw[cpdlink] (tI) -- (I);
  \draw[cpdlink] (tG) -- (G);
  \draw[cpdlink] (tS) -- (S);
  \draw[cpdlink] (tL) -- (L);
\end{tikzpicture}
\caption{The student network with its conditional probability distributions
(numbers from the lecture, after Koller \& Friedman). Each row of each table
is a probability distribution over the child, for one combination of parent
values.}
\label{fig:student}
\end{figure}

\subsection{How to read a CPD}

Variables without parents ($D$ and $I$) get an ordinary, unconditional
distribution. $P(I) = (0.7, 0.3)$ says that a random student is highly
motivated with probability $P(i^1) = 0.3$. Variables with parents get one
distribution \emph{per combination of parent values}. In $P(S \given I)$, the
row for $i^1$ says that a motivated student gets a low SAT score with
probability $P(s^0 \given i^1) = 0.2$ and a high one with probability $0.8$.
$P(G \given I, D)$ has four rows because its two parents have
$2 \times 2 = 4$ value combinations.

\begin{notebox}[title={Rows sum to one, columns do not}, unbreakable]
Every \emph{row} of a CPD is a distribution over the child and must sum to
one: $0.3 + 0.4 + 0.3 = 1$ for $(i^0, d^0)$. A \emph{column}, like the $g^1$
column $(0.3, 0.05, 0.9, 0.5)$, is a list of probabilities of the same child
value under different conditions and has no reason to sum to anything in
particular.
\end{notebox}

In the rest of the chapter we will see that this little picture is a lot more
than a sketch. Its structure supports different patterns of reasoning, it
captures all the independencies of the story, it is a compact representation
of the full joint distribution over all five variables --- precisely
\emph{because} of those independencies --- and therefore it can answer any
probabilistic query.

% ----------------------------------------------------------
\section{Reasoning Patterns}
\label{sec:reasoning-patterns}
% ----------------------------------------------------------

Remember the probabilistic query $P(\bX \given \be)$ from
\cref{sec:enumeration}: given the observed values $\be$ of the evidence
variables, what is the distribution over the query variables? Let us ask the
student network a sequence of such queries about a student called Joe and
watch how the answers move. For now, simply take the numbers as given; they
all follow from the CPDs in \cref{fig:student} (we compute two of them by hand
in \cref{sec:bn-event-probs}, and \cref{ch:inference} shows how to compute
such queries efficiently). \Cref{tab:joe} collects them.

\begin{table}[htbp]
\centering
\small
\begin{tabular}{@{}llll@{}}
\toprule
\textbf{What we know about Joe} & \textbf{Query} & \textbf{Answer} & \textbf{Pattern} \\
\midrule
nothing                         & $P(l^1)$                   & $0.502$ & --- \\
not motivated                   & $P(l^1 \given i^0)$        & $0.389$ & causal \\
not motivated, easy course      & $P(l^1 \given i^0, d^0)$   & $0.513$ & causal \\
\midrule
nothing                         & $P(i^1)$                   & $0.3$   & --- \\
grade C                         & $P(i^1 \given g^3)$        & $0.079$ & evidential \\
grade C                         & $P(d^1 \given g^3)$        & $0.63$  & evidential \\
\midrule
grade C, high SAT               & $P(i^1 \given g^3, s^1)$   & $0.578$ & intercausal \\
grade C, high SAT               & $P(d^1 \given g^3, s^1)$   & $0.76$  & intercausal (explaining away) \\
\bottomrule
\end{tabular}
\caption{Queries about Joe in the student network (rounded; all values
re-computed from the CPDs). For comparison, the priors are $P(d^1) = 0.4$ and
$P(i^1) = 0.3$.}
\label{tab:joe}
\end{table}

\paragraph{Causal reasoning (prediction)} Knowing nothing about Joe, his chance
of a strong letter is about $0.502$. When we learn that he is not very
motivated ($i^0$), it drops to $0.389$: unmotivated students tend to get worse
grades, and worse grades mean weaker letters. When we \emph{also} learn that
the course was easy ($d^0$), it rises again to $0.513$, because an easy course
makes a good grade likely despite his low motivation. Reasoning from causes to
effects --- ``downwards'' in the network --- is called \emph{causal reasoning}
or \emph{prediction}.

\paragraph{Evidential reasoning (explanation)} The prior probability that Joe
is highly motivated is $P(i^1) = 0.3$, read directly from the model. Now we
learn that he got a C ($g^3$). His probability of being motivated drops to
$0.079$. At the same time, the probability that the course was hard rises from
$0.4$ to $0.63$: a C is also evidence for a hard course. Reasoning from
observed effects to their possible causes --- ``upwards'' --- is called
\emph{evidential reasoning} or \emph{explanation}.

\paragraph{Intercausal reasoning} Besides the C, Joe now hands in his SAT
score, and surprisingly it is high ($s^1$). The probability that he is
motivated jumps from $0.079$ to $0.578$. Why does the SAT outweigh the poor
grade? Look at the CPDs: unmotivated students almost never get high SAT scores
($P(s^1 \given i^0) = 0.05$), while motivated students can still get a C
(especially in hard courses). But something else happens too: $P(d^1 \given
g^3, s^1) = 0.76$, up from $0.63$. The SAT score tells us about motivation,
and motivation --- \emph{together with the known grade} --- tells us about
difficulty. One cause of the grade (interest) gives us information about
another cause of the grade (difficulty). A poor grade can be due to low
motivation or to a hard course; once we learn that Joe is probably motivated,
the hard course becomes the more plausible explanation. This is
\emph{intercausal reasoning}; its best-known special case is called
\emph{explaining away}.

\begin{figure}[htbp]
\centering
\begin{tikzpicture}[
    flow/.style={-{Stealth[length=2.4mm]}, very thick, dashed, figamber},
    lab/.style={font=\small, align=center}]
  % (a) causal
  \begin{scope}
    \node[bnsmall, evid] (D) at (-0.8, 0)   {D};
    \node[bnsmall, evid] (I) at ( 0.8, 0)   {I};
    \node[bnsmall]       (G) at ( 0,  -1.3) {G};
    \node[bnsmall, qry]  (L) at ( 0,  -2.6) {L};
    \draw[edge] (D) -- (G); \draw[edge] (I) -- (G); \draw[edge] (G) -- (L);
    \draw[flow] (1.45,-0.2) to[bend left=25] (0.65,-2.6);
    \node[lab] at (0.2,-3.55) {(a) causal\\ $P(l^1 \given i^0, d^0)$};
  \end{scope}
  % (b) evidential
  \begin{scope}[xshift=4.6cm]
    \node[bnsmall, qry]  (D) at (-0.8, 0)   {D};
    \node[bnsmall, qry]  (I) at ( 0.8, 0)   {I};
    \node[bnsmall, evid] (G) at ( 0,  -1.3) {G};
    \node[bnsmall]       (L) at ( 0,  -2.6) {L};
    \draw[edge] (D) -- (G); \draw[edge] (I) -- (G); \draw[edge] (G) -- (L);
    \draw[flow] (0.75,-1.6) to[bend right=35] (1.45,0.2);
    \node[lab] at (0.2,-3.55) {(b) evidential\\ $P(i^1 \given g^3)$, $P(d^1 \given g^3)$};
  \end{scope}
  % (c) intercausal
  \begin{scope}[xshift=9.4cm]
    \node[bnsmall, qry]  (D) at (-0.8, 0)   {D};
    \node[bnsmall]       (I) at ( 0.8, 0)   {I};
    \node[bnsmall, evid] (G) at ( 0,  -1.3) {G};
    \node[bnsmall, evid] (S) at ( 1.9, -1.3) {S};
    \draw[edge] (D) -- (G); \draw[edge] (I) -- (G); \draw[edge] (I) -- (S);
    \draw[flow] (S.north) to[bend right=30] (I.east);
    \draw[flow] (I.south) .. controls (0.5,-2.6) and (-0.5,-2.6) .. (D.south);
    \node[lab] at (0.4,-3.55) {(c) intercausal\\ $P(d^1 \given g^3, s^1)$};
  \end{scope}
\end{tikzpicture}
\caption{The three reasoning patterns. Grey nodes are observed, red-framed
nodes are queried, and the dashed arrows show the direction in which
information flows. In (c), information about $I$ reaches $D$ only through the
\emph{observed} common effect $G$.}
\label{fig:reasoning-patterns}
\end{figure}

\begin{intbox}[unbreakable]
The textbook picture of explaining away: your lawn is wet in the morning. Rain
or the sprinkler? Both are plausible. Then you see that the sprinkler was on.
The wet lawn is now \emph{explained}, so rain becomes \emph{less} likely,
although rain and sprinkler have nothing to do with each other. The two causes
were independent \emph{until} we observed their common effect. The student
example is the mirror image: evidence makes one explanation of the C (low
motivation) \emph{less} likely, so the other one (a hard course) becomes
\emph{more} likely. Same mechanism: two causes of an observed effect compete
to explain it. This is the ``independent, but dependent given $Z$'' case from
\cref{sec:cond-indep}.
\end{intbox}

The remarkable thing is that we did not need a special algorithm for each
pattern. All of them --- including explaining away --- fall out of one
uniform inference procedure applied to the same model. A BN reasoner derives
\emph{every} consequence of \emph{any} combination of observations. Nor do we
need a notion of cause and effect for this to work: the arrows only describe
how the model is \emph{built}; information flows along them in both
directions. In principle, dependencies can even be modelled in any direction,
as long as the numbers are right --- although causal models are usually much
simpler (\cref{sec:building}).

% ----------------------------------------------------------
\section{Semantics I: The Independencies Encoded by the Graph}
\label{sec:bn-indep}
% ----------------------------------------------------------

We claimed that a BN is a model of the full joint distribution, and we know
from \cref{ch:prob} that independencies are what make a compact
representation possible. So let us first figure out which independencies the
student graph expresses. Recall the definition:
$(\bX \indep \bY \given \bZ)$ means $P(\bx \given \by, \bz) = P(\bx \given \bz)$
for all values, equivalently $P(\bX, \bY \given \bZ) = P(\bX \given \bZ)\,
P(\bY \given \bZ)$ --- the joint over $\bX, \bY, \bZ$ factorises.

\subsection{An intuitive analysis}

\paragraph{Letter} According to the model, the letter depends \emph{only} on
the grade. If we know the grade and the professor's habits (the CPD of $L$),
we can predict the letter, and learning anything else --- motivation, SAT,
difficulty --- changes nothing:
$(L \indep D, I, S \given G)$. But $L$ is \emph{not} independent of the other
variables if the grade is unknown: then motivation and difficulty do tell us
something about the letter, via the grade (that was causal reasoning). So
$(L \indep D, I, S \given G)$ but $(L \nindep D, I, S)$.

\paragraph{SAT} Motivation is the only direct influence on the SAT score. Once
we know Joe's motivation, his grade, the course difficulty or the letter tell
us nothing more about his SAT: $(S \indep D, G, L \given I)$. Without knowing
$I$, however, a good grade or letter does make a high SAT more likely:
$(S \nindep G, L)$.

\paragraph{Grade} The grade depends directly only on interest and difficulty.
Is it therefore independent of everything else given those two,
$(G \indep L, S \given I, D)$? \textbf{No!} Consider a strong student in a
hard course ($i^1, d^1$). The chance of an A is medium: $P(g^1 \given i^1, d^1)
= 0.5$. If we now learn that the student got a strong letter ($l^1$), our
belief in an A goes up: $P(g^1 \given i^1, d^1, l^1) \approx 0.712$. The
letter is an \emph{effect} of the grade, so observing it is evidence about the
grade, even when both parents are known: $(G \nindep L \given I, D)$.

What about $G$ and $S$? If we know motivation and difficulty, we have certain
expectations about the grade. Learning the SAT score could only change them by
changing our belief about motivation --- but motivation is already known. So
$(G \indep S \given I, D)$.

\paragraph{Difficulty and interest} These have no parents. Each of them can be
influenced by information about its descendants (evidential or intercausal
reasoning). But with no information about anything else, $I$ tells us nothing
about $D$, and vice versa; by the same argument $S$ tells us nothing about
$D$: $(D \indep I, S)$.

\subsection{The general pattern: local independencies}

All of these findings follow one rule. A variable's parents ``shield'' it from
the influence of everything further up in the network. Given its parents, a
variable is independent of all variables that are \emph{not its descendants}.
Information about its descendants, however, can still change our belief about
it, through evidential reasoning.

\begin{defbox}[title={Bayesian network structure (part I of the definition)}]
A \emph{Bayesian network structure} $\mathcal{G}$ is a directed acyclic graph
(DAG) whose nodes represent random variables $\Xall = \{X_1, \ldots, X_n\}$.
The graph encodes the following \emph{local independencies}:
\[
  (X_i \indep \NonDesc(X_i) \given \Pa(X_i)) \qquad\text{for all } X_i \in \Xall,
\]
where $\Pa(X_i)$ are the parents of $X_i$ in $\mathcal{G}$ and
$\NonDesc(X_i)$ are the variables that are not descendants of $X_i$.
In words: \emph{every variable is conditionally independent of its
non-descendants, given its parents.}
\end{defbox}

\Cref{fig:local-markov} illustrates the definition for $G$, and
\cref{tab:local-indep} lists the local independencies of all five variables.
They match the intuitive analysis exactly. All of them were also verified
numerically against the joint distribution defined by the CPDs.

\begin{figure}[htbp]
\centering
\begin{tikzpicture}
  \node[bn, fill=figblue!22] (D) at (-1.9, 0)   {Difficulty};
  \node[bn, fill=figblue!22] (I) at ( 1.9, 0)   {Interest};
  \node[bn, fill=accent!35, very thick] (G) at (-0.7,-1.7) {Grade};
  \node[bn, fill=figgreen!20] (S) at ( 3.1,-1.7) {SAT};
  \node[bn, fill=figred!12, dashed] (L) at (-0.7,-3.4) {Letter};
  \draw[edge] (D) -- (G); \draw[edge] (I) -- (G);
  \draw[edge] (I) -- (S); \draw[edge] (G) -- (L);
  % legend
  \begin{scope}[shift={(5.6,-0.2)}, font=\small]
    \node[draw=accent, fill=figblue!22, minimum size=3mm, inner sep=0pt] at (0,0) {};
    \node[anchor=west] at (0.25,0) {parents $\Pa(G) = \{D, I\}$};
    \node[draw=accent, fill=figgreen!20, minimum size=3mm, inner sep=0pt] at (0,-0.6) {};
    \node[anchor=west] at (0.25,-0.6) {other non-descendants: $S$};
    \node[draw=accent, dashed, fill=figred!12, minimum size=3mm, inner sep=0pt] at (0,-1.2) {};
    \node[anchor=west] at (0.25,-1.2) {descendants: $L$};
    \node[anchor=west, align=left] at (-0.15,-2.1)
      {$\Rightarrow\ (G \indep S \given D, I)$, \ but\ $(G \nindep L \given D, I)$};
  \end{scope}
\end{tikzpicture}
\caption{The local independence of $G$. Given its parents, $G$ is independent
of its other non-descendants, but not of its descendant $L$.}
\label{fig:local-markov}
\end{figure}

\begin{table}[htbp]
\centering
\small
\begin{tabular}{@{}lllll@{}}
\toprule
\textbf{Variable} & $\Pa$ & \textbf{Descendants} & \textbf{Non-desc.\ (besides parents)} & \textbf{Local independence} \\
\midrule
$D$ & ---       & $G, L$    & $I, S$    & $(D \indep I, S)$ \\
$I$ & ---       & $G, S, L$ & $D$       & $(I \indep D)$ \\
$G$ & $D, I$    & $L$       & $S$       & $(G \indep S \given D, I)$ \\
$S$ & $I$       & ---       & $D, G, L$ & $(S \indep D, G, L \given I)$ \\
$L$ & $G$       & ---       & $D, I, S$ & $(L \indep D, I, S \given G)$ \\
\bottomrule
\end{tabular}
\caption{Local independencies of the student network. Conditioning on the
parents is included in each statement; for a root node there is nothing to
condition on.}
\label{tab:local-indep}
\end{table}

\begin{notebox}[title={Two things to watch out for}]
\begin{itemize}
  \item \textbf{Descendants are excluded for a reason.} Observing a child (or
    grandchild) is evidence about its ancestors. This is exactly why
    $(G \nindep L \given I, D)$. A common mistake is to say ``a node is
    independent of everything given its parents'' --- it is independent only
    of its \emph{non-descendants}.
  \item \textbf{The list is not complete.} The local independencies are the
    ones the definition guarantees, but a graph usually implies
    \emph{additional} ones that follow from them. Finding all of them requires
    a full independence analysis of the graph (``d-separation''), which the
    lecture lists under advanced topics.
\end{itemize}
\end{notebox}

% ----------------------------------------------------------
\section{Semantics II: A Factorised Joint Distribution}
\label{sec:factorisation}
% ----------------------------------------------------------

How do the CPDs relate to the overall joint distribution $P$ over $\Xall$?
The answer: if the graph correctly represents the conditional independencies
of $P$, then $P$ can be rebuilt as a product of the local CPDs. The argument is
short and worth knowing.

\paragraph{Sketch of proof} Order the variables $X_1, \ldots, X_n$
\emph{topologically} with respect to the graph: every variable comes after all
of its parents (``top to bottom''; for the student network, e.g.\
$D, I, G, S, L$). Such an ordering always exists in an acyclic graph. Now apply
the chain rule in this order:
\[
P(X_1, \ldots, X_n) = P(X_1)\, P(X_2 \given X_1)\, P(X_3 \given X_1, X_2) \cdots
= \prod_{i=1}^{n} P(X_i \given X_1, \ldots, X_{i-1}).
\]
Because of the topological ordering, the predecessors $\{X_1, \ldots,
X_{i-1}\}$ contain all parents of $X_i$, possibly some other non-descendants,
and \emph{no} descendants. By the local independence $(X_i \indep
\NonDesc(X_i) \given \Pa(X_i))$, everything except the parents can be dropped
from each factor:
\[
P(X_i \given X_1, \ldots, X_{i-1}) = P(X_i \given \Pa(X_i)).
\]

\begin{thmbox}[title={Chain rule for Bayesian networks}, unbreakable]
\[
P(X_1, \ldots, X_n) = \prod_{i=1}^{n} P(X_i \given X_1, \ldots, X_{i-1})
                    = \prod_{i=1}^{n} P(X_i \given \Pa(X_i)).
\]
The CPDs jointly define a full probability distribution over $\Xall$; every
entry of the joint is a product of one entry from each CPD.
\end{thmbox}

\begin{defbox}[title={Bayesian network (part II of the definition)}]
A \emph{Bayesian network} $\mathcal{B}$ is a BN structure $\mathcal{G}$ in
which every node $X_i$ is associated with a CPD $P(X_i \given \Pa(X_i))$. The
network represents the joint distribution
\[
P(X_1, \ldots, X_n) = \prod_{i=1}^{n} P(X_i \given \Pa(X_i)).
\]
We say that $P$ \emph{factorises over} $\mathcal{G}$. A BN is a
\emph{factorised representation} of the joint distribution, and each CPD
$P(X_i \given \Pa(X_i))$ is called a \emph{factor}.
\end{defbox}

\begin{intbox}[unbreakable]
A BN can be read in two ways, and they describe the same object. Read as a
graph, it is a list of independence statements (``given its parents, each
variable ignores its non-descendants''). Read as a formula, it is a recipe for
the joint distribution (``multiply one number from each table''). The proof
above is the bridge: the independencies are exactly what allows the chain rule
to collapse into the short product.
\end{intbox}

For the student network, the factorisation reads
\[
P(D, I, G, S, L) = P(D)\, P(I)\, P(G \given D, I)\, P(S \given I)\, P(L \given G).
\]

\begin{exbox}[title={One entry of the joint}, unbreakable]
What is the probability that a random student is not very motivated, takes an
easy AI course, gets a B in it, has a low SAT score --- and still gets a
strong letter? Pick the matching entry from each CPD in \cref{fig:student}:
\begin{align*}
P(d^0, i^0, g^2, s^0, l^1)
&= P(d^0)\, P(i^0)\, P(g^2 \given d^0, i^0)\, P(s^0 \given i^0)\, P(l^1 \given g^2) \\
&= 0.6 \times 0.7 \times 0.4 \times 0.95 \times 0.6 = 0.09576 \approx 9.6\,\%.
\end{align*}
\end{exbox}

\subsection{Computing event probabilities from a BN}
\label{sec:bn-event-probs}

Since a BN \emph{is} a full joint distribution, everything from
\cref{sec:enumeration} carries over: the probability of an event is the sum of
all matching joint entries, and each joint entry is a product of CPD entries.
Usually you do not even need to write out the whole joint --- summing in a
clever order is enough. Here are two of the numbers from \cref{tab:joe}.

\begin{exbox}[title={The prior probability of a strong letter, $P(l^1)$}]
First get the marginal of $G$ by summing over its parents (law of total
probability, using $P(d, i) = P(d)\,P(i)$ since $D \indep I$):
\begin{align*}
P(g^1) &= \textstyle\sum_{i,d} P(i)\,P(d)\,P(g^1 \given i, d) \\
 &= 0.7\!\cdot\!0.6\!\cdot\!0.3 + 0.7\!\cdot\!0.4\!\cdot\!0.05
 + 0.3\!\cdot\!0.6\!\cdot\!0.9 + 0.3\!\cdot\!0.4\!\cdot\!0.5 = 0.362, \\
P(g^2) &= 0.42\!\cdot\!0.4 + 0.28\!\cdot\!0.25 + 0.18\!\cdot\!0.08
 + 0.12\!\cdot\!0.3 = 0.2884, \\
P(g^3) &= 1 - 0.362 - 0.2884 = 0.3496 .
\end{align*}
Then sum over the grade:
\[
P(l^1) = \sum_{g} P(g)\, P(l^1 \given g)
= 0.362 \cdot 0.9 + 0.2884 \cdot 0.6 + 0.3496 \cdot 0.01 \approx 0.502 .
\]
\end{exbox}

\begin{exbox}[title={Evidential reasoning by hand, $P(i^1 \given g^3)$}]
Bayes' rule with $P(g^3 \given i) = \sum_d P(d)\,P(g^3 \given i, d)$:
\begin{align*}
P(g^3 \given i^1) &= 0.6 \cdot 0.02 + 0.4 \cdot 0.2 = 0.092, \qquad
P(g^3 \given i^0) = 0.6 \cdot 0.3 + 0.4 \cdot 0.7 = 0.46, \\
P(i^1 \given g^3) &= \frac{P(g^3 \given i^1)\, P(i^1)}
  {P(g^3 \given i^1)\, P(i^1) + P(g^3 \given i^0)\, P(i^0)}
 = \frac{0.092 \cdot 0.3}{0.092 \cdot 0.3 + 0.46 \cdot 0.7}
 = \frac{0.0276}{0.3496} \approx 0.079 .
\end{align*}
Note that the denominator is again $P(g^3) = 0.3496$.
\end{exbox}

Doing this by hand gets tedious quickly, and the order in which we sum things
matters a lot for the effort. Organising these computations efficiently is
the topic of \cref{ch:inference}.

% ----------------------------------------------------------
\section{Compactness}
\label{sec:compactness}
% ----------------------------------------------------------

Recall the cost of the explicit joint. For $N$ Boolean variables the full
joint has $2^N$ entries, and specifying it takes $2^N - 1$ \emph{independent
parameters} --- the last entry is redundant because everything must sum to
one. Even for moderately sized worlds, a table like that is out of the
question.

Now consider a \emph{sparsely connected} BN over the same $N$ Boolean
variables, in which every variable has at most $k \ll N$ parents. Each row of a
CPD needs one number: the probability of $x_i^0$; the probability of $x_i^1$
is one minus that. A variable with $k$ parents has $2^k$ rows (one per parent
configuration), so its CPD needs $2^k$ numbers, and all CPDs together need at
most $N \cdot 2^k$ numbers. For fixed $k$ that is \emph{linear} in $N$: an
exponential reduction from $\bigO{2^N}$ to $\bigO{N}$.

\begin{thmbox}[title={Counting independent parameters}]
For variables with arbitrary finite domains,
\[
\#\text{params}(\text{full joint}) = \prod_{i=1}^{n} \abs{\Val(X_i)} - 1,
\]
\[
\#\text{params}(\text{BN}) = \sum_{i=1}^{n}
  \Big(\abs{\Val(X_i)} - 1\Big) \prod_{Y \in \Pa(X_i)} \abs{\Val(Y)} .
\]
Each CPD contributes (number of rows) $\times$ (number of child values $- 1$).
For Boolean variables with at most $k$ parents: $\leq N \cdot 2^k$.
\end{thmbox}

\begin{exbox}[title={Counting parameters in the examples}]
\textbf{Student network.} There are $2 \cdot 2 \cdot 3 \cdot 2 \cdot 2 = 48$
atomic events, so the explicit joint needs $47$ parameters. The BN needs
\[
\underbrace{1}_{D} + \underbrace{1}_{I} + \underbrace{4 \cdot 2}_{G} +
\underbrace{3 \cdot 1}_{L} + \underbrace{2 \cdot 1}_{S} = 15 .
\]
\textbf{Ten independent coins.} $2^{10} - 1 = 1023$ parameters for the
explicit joint; as a BN with no edges, $10 \times 1 = 10$.\par
\textbf{A realistic diagnosis model.} A technical system with 50 Boolean
component variables (ok/defective), each depending directly on at most 5
others: $2^{50} - 1 \approx 10^{15}$ parameters for the joint, but at most
$50 \times 2^5 = 1600$ for the BN.\par
\textbf{The diagnosis world of \cref{fig:diagnosis-net}.} $64 - 1 = 63$ for the
joint. As a BN: $\vn{Season}$ has 4 values (3 parameters), $\vn{Flu}$ and
$\vn{Hayfever}$ each have 4 rows (4 each), $\vn{MusclePain}$ has 2 rows (2),
$\vn{Congestion}$ has 4 rows (4): $3 + 4 + 4 + 2 + 4 = 17$.
\end{exbox}

\begin{notebox}[title={Subtract one per row, not one per table}, unbreakable]
In a CPD, \emph{every row} is its own distribution, so every row has one
redundant entry. A CPD for a three-valued child with two Boolean parents has
$4 \times 3 = 12$ entries but $4 \times 2 = 8$ independent parameters --- not
$11$. For the full joint, which is a single distribution, you subtract one in
total.
\end{notebox}

% ----------------------------------------------------------
\section{Building a Bayesian Network}
\label{sec:building}
% ----------------------------------------------------------

Constructing a BN means doing three things:
\begin{enumerate}
  \item \textbf{Choose the variables and their domains.} This is always the
    system designer's job.
  \item \textbf{Specify the structure.} Either by hand, with the designer and a
    domain expert working together (\emph{knowledge engineering}), or
    automatically by learning it from example situations (\emph{machine
    learning}, \cref{ch:learning}).
  \item \textbf{Specify the CPDs.} Usually learned from data; in realistic
    models there are far too many numbers to elicit by hand.
\end{enumerate}

One piece of advice matters a lot: \textbf{model in the direction of
causality}. Causal models are usually simpler, sparser and more intuitive than
models built ``backwards''. The slides leave it as an exercise to check this
for the student network; here is the solution.

\begin{exbox}[title={The student network, built backwards}]
Suppose we insist on the anti-causal ordering $L, G, S, D, I$ and, for each
variable, choose as parents the smallest set of its predecessors that makes it
independent of the remaining predecessors (so that the local independencies
hold). Using the true distribution of the student network:
\begin{itemize}
  \item $L$: no predecessors, no parents.
  \item $G$: depends on $L$, so $\Pa(G) = \{L\}$.
  \item $S$: given $G$ it is independent of $L$, so $\Pa(S) = \{G\}$.
  \item $D$: independent of $L$ given $G, S$ --- but \emph{not} independent of
    $S$ given $G$ (explaining away: given the grade, $S$ tells us about $I$,
    which tells us about $D$). So $\Pa(D) = \{G, S\}$.
  \item $I$: independent of $L$ given the rest, but depends on $G$, $S$ and
    $D$: $\Pa(I) = \{G, S, D\}$.
\end{itemize}
The result (\cref{fig:anticausal}) has 7 edges instead of 4 and needs
$1 + 2 \cdot 2 + 3 \cdot 1 + 6 \cdot 1 + 12 \cdot 1 = 26$ parameters instead of
15. It represents the same distribution, but it is bigger and much harder to
interpret --- try to explain to a domain expert what $P(I \given G, S, D)$
means! (All independence claims were checked numerically.)
\end{exbox}

\begin{figure}[htbp]
\centering
\begin{tikzpicture}
  % causal
  \begin{scope}
    \node[bnsmall] (D) at (-0.9, 0)   {D};
    \node[bnsmall] (I) at ( 0.9, 0)   {I};
    \node[bnsmall] (G) at ( 0,  -1.3) {G};
    \node[bnsmall] (S) at ( 1.8, -1.3) {S};
    \node[bnsmall] (L) at ( 0,  -2.6) {L};
    \draw[edge] (D) -- (G); \draw[edge] (I) -- (G);
    \draw[edge] (I) -- (S); \draw[edge] (G) -- (L);
    \node[font=\small, align=center] at (0.45,-3.5) {causal order\\ 4 edges, 15 parameters};
  \end{scope}
  % anti-causal
  \begin{scope}[xshift=4.3cm, yshift=-1.3cm]
    \node[bnsmall] (L) at (0,   0) {L};
    \node[bnsmall] (G) at (1.6, 0) {G};
    \node[bnsmall] (S) at (3.2, 0) {S};
    \node[bnsmall] (D) at (4.8, 0) {D};
    \node[bnsmall] (I) at (6.4, 0) {I};
    \draw[edge] (L) -- (G);
    \draw[edge] (G) -- (S);
    \draw[edge] (S) -- (D);
    \draw[edge] (D) -- (I);
    \draw[edge] (G) to[bend left=40] (D);
    \draw[edge] (G) to[bend left=50] (I);
    \draw[edge] (S) to[bend right=40] (I);
    \node[font=\small, align=center] at (3.2,-2.2) {anti-causal order $L, G, S, D, I$\\ 7 edges, 26 parameters};
  \end{scope}
\end{tikzpicture}
\caption{The same distribution with two different structures. Building the
network against the direction of causality forces extra edges.}
\label{fig:anticausal}
\end{figure}

\subsection{Two larger examples}

\Cref{fig:car} shows a simplified BN for car diagnosis (originally from
lecture slides by Stuart Russell). The colours show the different roles that
variables play in a diagnostic model: the \emph{problem} to be explained (the
car won't start), \emph{tests} whose values can be obtained (lights, gauges,
dipstick), possible \emph{diagnoses} that are not directly observable (dead
battery, no oil, broken starter, \ldots), and \emph{auxiliary} variables. The
auxiliary nodes, such as ``battery flat'' or ``no charging'', do not
correspond to anything we would test or repair. They exist to make the model
more structured and compact: many causes act on the tests \emph{through} a
flat battery, so they only need one shared intermediate node instead of
separate edges to every test.

\begin{figure}[htbp]
\centering
\begin{tikzpicture}[
    car/.style={ellipse, draw=accent, thick, font=\scriptsize, align=center,
                minimum width=17mm, minimum height=8.5mm, inner sep=0.5pt},
    test/.style={car, fill=figgreen!22},
    diag/.style={car, fill=figamber!30, very thick},
    aux/.style={car, fill=gray!22},
    prob/.style={car, fill=figred!30, very thick},
    cedge/.style={-{Stealth[length=2mm]}, semithick, accent!85!black}]
  \node[test] (age)  at (-5.6, 0)    {battery\\age};
  \node[diag] (alt)  at (-2.4, 0)    {alternator\\broken};
  \node[diag] (fan)  at ( 0.0, 0)    {fanbelt\\broken};
  \node[diag] (dead) at (-4.6,-1.45) {battery\\dead};
  \node[aux]  (chg)  at (-1.2,-1.45) {no\\charging};
  \node[test] (mtr)  at (-6.0,-2.9)  {battery\\meter};
  \node[aux]  (flat) at (-3.2,-2.9)  {battery\\flat};
  \node[diag] (oil)  at (-0.9,-2.9)  {no oil};
  \node[diag] (gas)  at ( 1.3,-2.9)  {no gas};
  \node[diag] (fuel) at ( 3.5,-2.9)  {fuel line\\blocked};
  \node[diag] (strt) at ( 5.7,-2.9)  {starter\\broken};
  \node[test] (lgt)  at (-5.4,-4.6)  {lights};
  \node[test] (oill) at (-3.0,-4.6)  {oil light};
  \node[test] (gg)   at (-0.8,-4.6)  {gas gauge};
  \node[prob] (wont) at ( 1.6,-4.6)  {car won't\\start};
  \node[test] (dip)  at ( 4.0,-4.6)  {dipstick};
  \foreach \a/\b in {age/dead, alt/chg, fan/chg, dead/mtr, dead/flat, chg/flat,
                     flat/lgt, flat/oill, flat/gg, flat/wont,
                     oil/oill, oil/wont, oil/dip, gas/gg, gas/wont,
                     fuel/wont, strt/wont}
    \draw[cedge] (\a) -- (\b);
  % legend
  \begin{scope}[shift={(-4.4,-5.75)}, font=\scriptsize]
    \foreach \st/\txt/\x/\y in {prob/problem to explain (observed)/0/0,
                                test/test (can be observed)/5.4/0,
                                diag/diagnosis (hidden)/0/-0.5,
                                aux/auxiliary (keeps the model compact)/5.4/-0.5} {
      \node[\st, minimum width=6mm, minimum height=3.2mm, inner sep=0pt] at (\x,\y) {};
      \node[anchor=west] at (\x+0.4,\y) {\txt};
    }
  \end{scope}
\end{tikzpicture}
\caption{A simplified Bayesian network for car diagnosis (after S.~Russell).
Arrows point from causes to effects; ``car won't start'' is the problem to be
explained.}
\label{fig:car}
\end{figure}

The second example is a real one: the ALARM network for monitoring patients
in an intensive care unit (Beinlich et al., 1989), with 37 variables. Despite
its size, every variable depends directly on only a few others, which is
exactly what keeps the model small enough to specify and to compute with.

\figureplaceholder{5.5cm}{The ALARM network for ICU monitoring}{A Bayesian
network with 37 variables for monitoring ventilated patients
(e.g.\ HYPOVOLEMIA, LVFAILURE, SAO2, HR, BP, VENTLUNG). Any clean rendering of
the full graph works.\par\smallskip
\textit{Search hint:} ``ALARM Bayesian network Beinlich 1989'' --- e.g.\ the
graph on the bnlearn repository page for the ALARM network.}{The ALARM network
(Beinlich et al., 1989): a realistic medical Bayesian network with 37
variables, each with only a few parents.}{fig:alarm}

% ----------------------------------------------------------
\section{Continuous Variables}
\label{sec:bn-continuous}
% ----------------------------------------------------------

So far all variables were discrete, which let us write every CPD
$P(X \given \Pa(X))$ as a table. But many systems involve continuous
quantities --- the position and velocity of an aircraft or a robot, and our
measurements of them --- which are best modelled by real-valued variables.
There are two ways to deal with them.

\paragraph{Solution 1: discretisation} Split the domain of the variable into a
fixed number of intervals and represent each value by the label of its
interval. The advantage is that all of our machinery applies unchanged. The
disadvantages are the information lost in the approximation and, above all,
complexity. For reasonably accurate robot tracking we might need 1000
intervals per position coordinate, i.e.\ $1000 \times 1000 = 10^6$ different
$(x, y)$ positions. Now imagine a CPD with that as a parent.

\paragraph{Solution 2: continuous variables with continuous CPDs} Allow
$\Val(X) \subseteq \R$ and model CPDs as density functions $p(X \given
\Pa(X))$. This raises two questions: how do we model the distribution of a
continuous variable, and how do we define $p(X \given \Pa(X))$ when the
parents are continuous, too? Take $X \to Y$ with both continuous. The CPD
$p(Y \given X)$ must contain one distribution $p(Y \given x)$ for
\emph{every} value $x$ --- infinitely many. We cannot list them.

The trick is to pick a \emph{parametric} family of densities and to make its
parameters a \emph{function} of the parent values. Then a few numbers describe
all the infinitely many conditional distributions at once. The most important
special case:

\begin{defbox}[title={Linear Gaussian model}]
Let $Y$ be a continuous variable with continuous parents $X_1, \ldots, X_k$.
$Y$ has a \emph{linear Gaussian model} if there are parameters
$\beta_0, \beta_1, \ldots, \beta_k$ and $\sigma^2$ such that
\[
p(Y \given x_1, \ldots, x_k) = \Normal(\beta_0 + \beta_1 x_1 + \cdots + \beta_k x_k;\ \sigma^2),
\]
or in vector notation $p(Y \given \bx) = \Normal(\beta_0 + \bm{\beta}\T \bx;\ \sigma^2)$.
A \emph{linear Gaussian Bayesian network} is a BN in which all variables are
continuous and all of them have linear Gaussian models.
\end{defbox}

So every distribution over a single variable is a Gaussian; the \emph{mean} of
the child is a linear function of the parents; and the \emph{variance} is
fixed --- it does not depend on the parents. For example, $p(Y \given x) =
\Normal(3.5x - 0.9;\ 1.0)$: whatever $x$ is, $Y$ is spread around $3.5x - 0.9$
with variance one. An equivalent and often more intuitive way to say this: $Y$
is a linear function of its parents plus Gaussian noise,
\[
y = \beta_0 + \beta_1 x_1 + \cdots + \beta_k x_k + \epsilon,
\qquad \epsilon \sim \Normal(0, \sigma^2).
\]
\Cref{fig:linear-gaussian} shows what this looks like for one parent.

\begin{figure}[htbp]
\centering
\begin{tikzpicture}[scale=1.15, font=\small]
  \draw[-{Stealth[length=2mm]}, gray] (0,0) -- (4.6,0) node[right] {$x$ (parent)};
  \draw[-{Stealth[length=2mm]}, gray] (0,0) -- (0,4.4) node[above] {$y$ (child)};
  % samples
  \foreach \p in {(0.32,1.11), (0.53,1.38), (0.86,1.22), (1.10,1.56), (1.28,1.52),
                  (1.53,1.87), (1.84,2.14), (2.06,2.11), (2.37,2.87), (2.58,2.54),
                  (2.93,3.51), (3.07,3.22), (3.29,3.42), (3.65,3.69), (3.80,3.27),
                  (4.08,3.65)}
    \fill[gray!60] \p circle (1.3pt);
  % mean line
  \draw[figblue, very thick] (0.1,0.68) -- (4.3,4.04)
    node[right, align=left] {mean\\ $\beta_0 + \beta_1 x$};
  % gaussian slices
  \foreach \xx in {1, 2.2, 3.4} {
    \pgfmathsetmacro{\ymean}{0.6 + 0.8*\xx}
    \draw[gray, dashed] (\xx, 0) -- (\xx, \ymean);
    \draw[figred, thick, domain=-1:1, samples=60, smooth]
      plot ({\xx + 0.75*exp(-(\x)^2/(2*0.09))}, {\ymean + \x});
  }
  \node[figred, anchor=west, align=left] at (0.25,4.0)
    {$p(Y \given x) = \Normal(\beta_0 + \beta_1 x;\ \sigma^2)$};
\end{tikzpicture}
\caption{A linear Gaussian CPD with one parent. For every parent value $x$ the
child is Gaussian (red curves); the mean moves along a straight line, while the
width stays the same everywhere. Grey dots are samples from the model.}
\label{fig:linear-gaussian}
\end{figure}

\begin{exbox}[title={A bicycle on a straight road}]
A cyclist rides along a straight road; once per second we model her position
$X^{(t)}$ (in m) and her velocity $V^{(t)}$ (in m/s). Ideally,
$X^{(t+1)} = X^{(t)} + V^{(t)}$: at metre 510.4 with a speed of 5\,m/s, she
should be at metre 515.4 one second later. But real motion is never perfectly
precise, so it is more realistic to say that the next position is normally
distributed around that value:
\[
p(X^{(t+1)} \given x^{(t)}, v^{(t)}) = \Normal(x^{(t)} + v^{(t)};\ 0.8).
\]
This is a linear Gaussian model with $\beta_0 = 0$, $\beta_1 = \beta_2 = 1$ and
$\sigma^2 = 0.8$. (The slides say ``variance 0.8 meters''; strictly a variance
is measured in m$^2$ --- the standard deviation is $\sqrt{0.8} \approx
0.89$\,m.)
\end{exbox}

Gaussian distributions, linear dependencies and a fixed variance are fairly
severe restrictions. Some processes satisfy them and some do not. But they are
reasonably good models of many physical phenomena, and --- the real selling
point --- they support simple and elegant reasoning algorithms. We will meet
a concrete one for a restricted family of such models in the chapter on the
Kalman filter (\cref{ch:kalman}).

\paragraph{Beyond this chapter} The lecture explicitly leaves out a full
analysis of independence in BN graphs, algorithms that construct a graph for a
given set of independencies, and other kinds of local models (context-specific
CPDs, tree and rule CPDs, hybrid discrete--continuous models, logistic models,
\ldots).

% ----------------------------------------------------------
\section{Chapter Checklist}
% ----------------------------------------------------------

What you should take away from this chapter (following the lecture's own
list):
\begin{itemize}
  \item the probabilistic query, and the typical reasoning patterns: causal,
    evidential and intercausal (explaining away);
  \item the definition of a Bayesian network, in both parts: the graph with
    its local independencies, and the CPDs with the factorised distribution;
  \item which independencies a BN encodes ---
    $(X_i \indep \NonDesc(X_i) \given \Pa(X_i))$ --- and why descendants are
    excluded;
  \item the chain rule for Bayesian networks and its proof via a topological
    ordering;
  \item how to compute event probabilities from a BN;
  \item the notion of independent (non-redundant) parameters, and how a BN
    makes the representation of the full joint compact;
  \item the basic idea of linear Gaussian models (details follow with the
    Kalman filter).
\end{itemize}


% ==========================================================
%  CHAPTERS 4--8 — OUTLINES
%  Section lists follow the lecture titles in the course schedule
%  and the standard treatment in Koller & Friedman / Russell & Norvig.
%  Replace \planned{...} blocks with real text once the slides of the
%  corresponding units are available, and adapt the sections to them.
% ==========================================================

% ==========================================================
%  CHAPTER 4 — INFERENCE IN BAYESIAN NETWORKS   (VO 4--5)
% ==========================================================
\chapter{Inference in Bayesian Networks}
\label{ch:inference}

\Cref{ch:bn} showed that a Bayesian network \emph{is} a full joint
distribution, so every query can in principle be answered by enumeration ---
multiply CPD entries, sum over the hidden variables, normalise. In practice
that is exponential in the number of variables. This chapter is about doing it
smarter: first exactly, by organising the sums cleverly, and then
approximately, by \emph{sampling} from the network when exact inference is too
expensive.

\planned{Lecture units VO~4 (``Bayes Nets: Exact \& Approximate Inference
(1)'') and VO~5 (``Bayes Nets: Approximate Inference and MCMC''). The
section list below is the expected structure; adapt it to the slides.}

\section{Exact Inference by Enumeration in a Bayesian Network}
\planned{Recap of the probabilistic query; enumeration with the factorised
joint; pushing sums inwards; the student network as a worked example
(re-using the numbers of \cref{tab:joe}); cost of naive enumeration.}

\section{Variable Elimination}
\planned{Factors and factor operations (product, summing out, reduction by
evidence); the variable-elimination algorithm; worked example; the effect of
the elimination order on cost; intermediate factor sizes. Figure (TikZ, to be
drawn): a small elimination example showing which factors are created.}

\section{The Complexity of Exact Inference}
\planned{Exact inference is NP-hard in general; polytrees (singly connected
networks) allow inference in linear time; why loops make it hard; motivation
for approximate methods.}

\section{Approximate Inference by Stochastic Sampling}
\planned{The idea: estimate probabilities from frequencies in samples. Prior
(forward) sampling in topological order; rejection sampling for conditional
queries and its weakness with unlikely evidence; likelihood weighting;
consistency and convergence of the estimates.}

\section{Markov Chain Monte Carlo and Gibbs Sampling}
\planned{Markov chains over complete states; Gibbs sampling: resample one
variable at a time given its Markov blanket; stationary distribution; burn-in
and mixing; comparison with likelihood weighting. Figure (TikZ, to be drawn):
a chain of complete assignments where each step resamples one variable.}


\section{Chapter Checklist}
\planned{To be written together with the chapter.}


% ==========================================================
%  CHAPTER 5 — LEARNING BAYESIAN NETWORKS   (VO 6--8)
% ==========================================================
\chapter{Learning Bayesian Networks}
\label{ch:learning}

So far somebody handed us the graph and all the CPDs. In reality, the numbers
--- and often the structure --- have to come from data. This chapter addresses
the third of the complexity problems from \cref{ch:intro}, modelling
complexity: given a set of observations, find the parameters (and maybe the
structure) of a Bayesian network that explains them well. Two schools of
parameter estimation appear --- maximum likelihood and Bayesian --- followed by
the question of how to avoid overfitting and how to search over structures.

\planned{Lecture units VO~6 (``Learning (1): General Issues \& ML Parameter
Learning''), VO~7 (``Learning (2): Bayesian Parameter Learning'') and VO~8
(``Learning (3): Overfitting, Regularisation, Structure Learning''). Expected
structure below; adapt it to the slides.}

\section{The Learning Problem}
\planned{What is given and what is learned (parameters vs.\ structure);
complete vs.\ incomplete data; the i.i.d.\ assumption; goals of learning
(density estimation, prediction, knowledge discovery).}

\section{Maximum-Likelihood Parameter Learning}
\planned{The likelihood function; ML estimation for a single Bernoulli/
multinomial variable (``counting''); decomposition of the BN likelihood into
independent per-CPD problems; sufficient statistics; worked example with a
small data table; the zero-count problem.}

\section{Bayesian Parameter Learning}
\planned{Parameters as random variables; prior, posterior and predictive
distribution; Beta and Dirichlet priors; conjugacy; pseudo-counts and their
intuition; MAP vs.\ full Bayesian prediction; worked coin example. Figure
(pgfplots, to be drawn): Beta priors and how the posterior sharpens with
data.}


\section{Overfitting and Regularisation}
\planned{Why maximum likelihood overfits, especially for CPDs with many
parents and little data; priors as regularisers; the bias--variance view;
model complexity.}

\section{Structure Learning}
\planned{Score-based structure learning: likelihood score (and why it always
prefers more edges), penalised scores (BIC/MDL), Bayesian score; local search
over graphs (add/delete/reverse an edge); decomposability of scores.}

\section{Chapter Checklist}
\planned{To be written together with the chapter.}


% ==========================================================
%  CHAPTER 6 — TEMPORAL MODELS: DYNAMIC BAYESIAN NETWORKS   (VO 9)
% ==========================================================
\chapter[Temporal Models: Dynamic Bayesian Networks]{Temporal Models:\\ Dynamic Bayesian Networks}
\label{ch:dbn}

Many of the examples from \cref{ch:intro} --- tracking aircraft, driving a car,
following music --- are about a world that \emph{changes over time}. We
observe it step by step and want to keep our beliefs up to date. Dynamic
Bayesian networks extend the BN idea to such processes: the same small
``template'' of variables is repeated for every time step, with edges inside a
time slice and from one slice to the next.

\planned{Lecture unit VO~9 (``Dynamic Bayes Nets and State-Observation
Models''). Expected structure below; adapt it to the slides.}

\section{Template Variables and Time Slices}
\planned{Discrete time, template variables $X_i^{(t)}$ (see
\cref{tab:notation}); the Markov assumption; stationarity
(time-invariant transition model); the two-slice representation (2-TBN) and
unrolling it over $T$ steps.}

\section{State-Observation Models}
\planned{Hidden state and observations; transition model $P(\bX^{(t+1)} \given
\bX^{(t)})$ and sensor/observation model $P(\bY^{(t)} \given \bX^{(t)})$;
examples (robot localisation, the bicycle of \cref{sec:bn-continuous}).
Figure (TikZ, to be drawn): an unrolled state-observation chain.}

\section{Inference Tasks over Time}
\planned{Filtering (belief state), prediction, smoothing, most likely
explanation; the recursive forward update (predict, then condition on the new
observation); why naive unrolling becomes expensive.}

\section{Chapter Checklist}
\planned{To be written together with the chapter.}


% ==========================================================
%  CHAPTER 7 — HIDDEN MARKOV MODELS   (VO 10--11)
% ==========================================================
\chapter{Hidden Markov Models}
\label{ch:hmm}

A hidden Markov model is the simplest interesting state-observation model: a
single discrete hidden state that evolves as a Markov chain, and an
observation emitted at every step. Despite (or because of) this simplicity,
HMMs power speech recognition, gene finding and --- close to the lecturer's
research --- following a live music performance in a score. This chapter
covers the classic inference algorithms, learning the parameters from data,
and applications.

\planned{Lecture units VO~10 (``Hidden Markov Models (1): Inference'') and
VO~11 (``Hidden Markov Models (2): Learning and Applications''). Expected
structure below; adapt it to the slides.}

\section{The HMM as a Dynamic Bayesian Network}
\planned{States, observations, initial distribution $\pi$, transition matrix
$A$, emission probabilities $B$; the HMM as a special DBN; a small running
example. Figure (TikZ, to be drawn): HMM unrolled over time, plus the state
transition diagram.}

\section{Filtering and Likelihood: The Forward Algorithm}
\planned{Forward messages $\alpha_t$; recursion and normalisation; computing
the likelihood of an observation sequence; complexity $\bigO{T N^2}$
instead of $\bigO{N^T}$.}

\section{Smoothing: The Forward--Backward Algorithm}
\planned{Backward messages $\beta_t$; combining forward and backward messages
into posteriors over past states.}

\section{The Most Likely State Sequence: Viterbi}
\planned{Max instead of sum (compare the max-marginal of
\cref{sec:enumeration}); dynamic programming over the trellis; backtracking;
worked example. Figure (TikZ, to be drawn): Viterbi trellis with the best path
highlighted.}

\section{Learning HMMs: Baum--Welch}
\planned{Supervised (counting) vs.\ unsupervised learning; the EM idea;
expected counts from forward--backward; re-estimation formulas; local optima.}

\section{Applications}
\planned{Speech recognition; real-time music tracking / score following; other
applications mentioned in the lecture.}

\figureplaceholder{4.5cm}{Real-time score following}{An image of a score
follower in action: a live audio signal aligned to the sheet music, with the
current position highlighted (as in the lecture's video of real-time music
following).\par\smallskip
\textit{Search hint:} ``automatic score following live performance'' or
``real-time music tracking score alignment''; screenshots of the JKU
Computational Perception group's demos fit well.}{An HMM-based score follower
tracks a live performance in the printed score.}{fig:score-following}

\section{Chapter Checklist}
\planned{To be written together with the chapter.}


% ==========================================================
%  CHAPTER 8 — THE KALMAN FILTER   (video lecture)
% ==========================================================
\chapter[Continuous Linear Systems: The Kalman Filter]{Continuous Linear Systems:\\ The Kalman Filter}
\label{ch:kalman}

What an HMM does for a discrete hidden state, the Kalman filter does for a
continuous one --- provided the world is \emph{linear} and \emph{Gaussian}.
This is exactly the linear Gaussian setting previewed in
\cref{sec:bn-continuous}, and it is what makes the bicycle example tractable:
if the belief about the current state is Gaussian and all dependencies are
linear Gaussian, then the belief stays Gaussian forever, and filtering reduces
to updating a mean and a covariance.

\planned{Video lecture ``Modelling Continuous Linear Systems: The Kalman
Filter''. Expected structure below; adapt it to the video.}

\section{Linear Gaussian State-Space Models}
\planned{Continuous state and observation variables; linear Gaussian
transition and sensor models (connection to \cref{sec:bn-continuous});
the bicycle/robot tracking example.}

\section{The Predict--Update Cycle}
\planned{Prediction step (the belief spreads out); update step (conditioning
on the observation sharpens it); the Kalman gain as a weighting between
prediction and measurement; the 1-D case in detail, then the matrix form.
Figure (pgfplots, to be drawn): prior, prediction, measurement and posterior
Gaussians in 1-D.}

\section{Properties and Limitations}
\planned{Why Gaussians stay Gaussian; optimality under the model assumptions;
what goes wrong for non-linear or multi-modal problems; extensions if covered
in the video.}

\section{Chapter Checklist}
\planned{To be written together with the chapter.}


\end{document}
